% ============================================================================
% 2026-09-30 (agent L4): the 16:00 campaign (writeup/CAMPAIGN_16H_2026-09-29.md) as a
% record, one change at a time.  Every number below is a macro from
% generated/appendix_campaign_16h_numbers.tex and every table a file
% generated/appendix_campaign_16h_<table>.tex, both written by
%     python writeup/make_appendix_campaign_16h_tex.py
% from the campaign's CSVs (each table's source is named in a comment above it and
% inside the generated file); nothing is typed by hand.  Wording: sign(s) rule,
% straddle, last-30-min trade; no cluster or lab names in prose (the compute table
% labels them cluster A / B).  Input from main.tex after
% sections/appendix_close_compounding, so this continues Appendix D.
% ============================================================================
% AUTO-GENERATED by writeup/make_appendix_campaign_16h_tex.py -- do not edit.
% Every macro is read from the campaign's CSVs / summaries (sources in the table files and
% in the script); prose in sections/appendix_campaign_16h.tex quotes numbers only through these.
\newcommand{\campPOldBase}{34}
\newcommand{\campPNewBase}{22}
\newcommand{\campKeptLinBase}{17}
\newcommand{\campPOldLive}{244}
\newcommand{\campPNewLive}{232}
\newcommand{\campKeptLinLive}{136--138}
\newcommand{\campPOldAll}{640}
\newcommand{\campPNewAll}{628}
\newcommand{\campKeptLinAll}{360--378}
\newcommand{\campNSessionEdge}{12}
\newcommand{\campNSessionEdgeHalf}{6}
\newcommand{\campNLinear}{76}
\newcommand{\campNLinearMasked}{75}
\newcommand{\campMaskDropAgree}{75}
\newcommand{\campOLSN}{1}
\newcommand{\campOLSMaxRel}{\ensuremath{8\times10^{-15}}}
\newcommand{\campOLSMoved}{0}
\newcommand{\campRidgeN}{25}
\newcommand{\campRidgeMaxRel}{\ensuremath{4\times10^{-10}}}
\newcommand{\campRidgeMoved}{0}
\newcommand{\campLassoN}{25}
\newcommand{\campLassoMaxRel}{\ensuremath{9\times10^{-2}}}
\newcommand{\campLassoMoved}{7}
\newcommand{\campEnetN}{25}
\newcommand{\campEnetMaxRel}{\ensuremath{9\times10^{-2}}}
\newcommand{\campEnetMoved}{10}
\newcommand{\campMovedTol}{\ensuremath{1\times10^{-8}}}
\newcommand{\campPosChanged}{9}
\newcommand{\campTradeDays}{866}
\newcommand{\campMaxDSharpeLin}{0.06}
\newcommand{\campAttrN}{17}
\newcommand{\campAttrRepeat}{17}
\newcommand{\campAttrDesignN}{13}
\newcommand{\campAttrDesignLo}{\ensuremath{5.5\times10^{-4}}}
\newcommand{\campAttrDesignHi}{\ensuremath{1.7\times10^{-2}}}
\newcommand{\campAttrMachineN}{12}
\newcommand{\campAttrMachineLo}{\ensuremath{7.3\times10^{-4}}}
\newcommand{\campAttrMachineHi}{\ensuremath{8.1\times10^{-2}}}
\newcommand{\campLSTMDedupLo}{2.5}
\newcommand{\campLSTMDedupHi}{7.6}
\newcommand{\campRodChecks}{9}
\newcommand{\campRodBeforeMax}{\ensuremath{1.7\times10^{-2}}}
\newcommand{\campRodBeforeSameMin}{865}
\newcommand{\campRodDeck}{866}
\newcommand{\campRodAfterBit}{9}
\newcommand{\campRodRows}{1{,}469}
\newcommand{\campRodArms}{117}
\newcommand{\campRodMaskSame}{117}
\newcommand{\campKeptBase}{17}
\newcommand{\campKeptBaseLo}{17}
\newcommand{\campKeptBaseHi}{17}
\newcommand{\campPBase}{22}
\newcommand{\campKeptLive}{138}
\newcommand{\campKeptLiveLo}{136}
\newcommand{\campKeptLiveHi}{145}
\newcommand{\campPLive}{232}
\newcommand{\campKeptAll}{369}
\newcommand{\campKeptAllLo}{360}
\newcommand{\campKeptAllHi}{378}
\newcommand{\campPAll}{628}
\newcommand{\campKeptRefits}{1{,}469}
\newcommand{\campMaskTreeN}{36}
\newcommand{\campMaskTreeQBelow}{0}
\newcommand{\campMaskTreeQAbove}{5}
\newcommand{\campMaskTreeSBelow}{1}
\newcommand{\campMaskTreeSAbove}{0}
\newcommand{\campMaskMedAbsQ}{0.7}
\newcommand{\campMaskLSTMN}{3}
\newcommand{\campMaskLSTMQBelow}{1}
\newcommand{\campMaskLSTMQAbove}{0}
\newcommand{\campMaskLSTMSBelow}{0}
\newcommand{\campMaskLSTMSAbove}{0}
\newcommand{\campMaskLSTMAllQ}{\ensuremath{-}17.1}
\newcommand{\campMaskLSTMAllQCI}{[\ensuremath{-}33.2, \ensuremath{-}4.2]}
\newcommand{\campTOneTenN}{9}
\newcommand{\campTOneTenQBelow}{3}
\newcommand{\campTOneTenQAbove}{0}
\newcommand{\campTOneTenSBelow}{0}
\newcommand{\campTOneTenSAbove}{1}
\newcommand{\campTOneTenQLo}{\ensuremath{-}3.6}
\newcommand{\campTOneTenQHi}{\ensuremath{-}0.8}
\newcommand{\campRSOneTenN}{9}
\newcommand{\campRSOneTenQBelow}{1}
\newcommand{\campRSOneTenQAbove}{0}
\newcommand{\campRSOneTenSBelow}{0}
\newcommand{\campRSOneTenSAbove}{0}
\newcommand{\campRSOneTenQLo}{\ensuremath{-}2.2}
\newcommand{\campRSOneTenQHi}{\ensuremath{+}0.2}
\newcommand{\campRSVsTN}{9}
\newcommand{\campRSVsTQBelow}{0}
\newcommand{\campRSVsTQAbove}{6}
\newcommand{\campRSVsTSBelow}{0}
\newcommand{\campRSVsTSAbove}{0}
\newcommand{\campRSVsTQLo}{\ensuremath{+}3.6}
\newcommand{\campRSVsTQHi}{\ensuremath{+}17.9}
\newcommand{\campLSTMCadN}{3}
\newcommand{\campLSTMCadQBelow}{2}
\newcommand{\campLSTMCadQAbove}{0}
\newcommand{\campLSTMCadSBelow}{0}
\newcommand{\campLSTMCadSAbove}{0}
\newcommand{\campLSTMCadBase}{\ensuremath{-}1.9}
\newcommand{\campLSTMCadBaseCI}{[\ensuremath{-}3.5, \ensuremath{-}0.3]}
\newcommand{\campLSTMCadLive}{\ensuremath{-}2.5}
\newcommand{\campLSTMCadLiveCI}{[\ensuremath{-}7.8, \ensuremath{+}3.2]}
\newcommand{\campLSTMCadAll}{\ensuremath{-}10.1}
\newcommand{\campLSTMCadAllCI}{[\ensuremath{-}18.0, \ensuremath{-}2.0]}
\newcommand{\campRidgeTreeN}{36}
\newcommand{\campRidgeTreeQBelow}{0}
\newcommand{\campRidgeTreeQAbove}{18}
\newcommand{\campRidgeTreeSBelow}{0}
\newcommand{\campRidgeTreeSAbove}{0}
\newcommand{\campRidgeBestTree}{T1 LightGBM \texttt{live\_feasible}}
\newcommand{\campRidgeBestTreeQ}{\ensuremath{-}5.0}
\newcommand{\campRidgeBestTreeQCI}{[\ensuremath{-}11.5, \ensuremath{+}0.4]}
\newcommand{\campRidgeLSTMN}{3}
\newcommand{\campRidgeLSTMQBelow}{0}
\newcommand{\campRidgeLSTMQAbove}{3}
\newcommand{\campRidgeLSTMSBelow}{0}
\newcommand{\campRidgeLSTMSAbove}{0}
\newcommand{\campRidgeLSTMBase}{\ensuremath{+}6.8}
\newcommand{\campRidgeLSTMLive}{\ensuremath{+}18.6}
\newcommand{\campRidgeLSTMAll}{\ensuremath{+}30.0}
\newcommand{\campOptTrials}{50}
\newcommand{\campOptPoints}{1{,}469}
\newcommand{\campOptExch}{20}
\newcommand{\campOptLateLo}{27}
\newcommand{\campOptLateHi}{40}
\newcommand{\campOptGainTwentyFiveLo}{1.1}
\newcommand{\campOptGainTwentyFiveHi}{2.1}
\newcommand{\campOptGainTwentyFiveMedLo}{0.4}
\newcommand{\campOptGainTwentyFiveMedHi}{1.2}
\newcommand{\campOptGainFortyLo}{0.3}
\newcommand{\campOptGainFortyHi}{0.7}
\newcommand{\campOptGainFortyMedMax}{0.0}
\newcommand{\campOptPickDiffLo}{61}
\newcommand{\campOptPickDiffHi}{74}
\newcommand{\campOptShippedHi}{1.1}
\newcommand{\campBudN}{54}
\newcommand{\campBudQBelow}{0}
\newcommand{\campBudQAbove}{1}
\newcommand{\campBudSBelow}{1}
\newcommand{\campBudSAbove}{0}
\newcommand{\campTPN}{27}
\newcommand{\campTPQBelow}{8}
\newcommand{\campTPQAbove}{0}
\newcommand{\campTPSBelow}{0}
\newcommand{\campTPSAbove}{0}
\newcommand{\campTPBelowWho}{LightGBM \texttt{baseline} (TUNE\_PER 5, 25); random forest \texttt{live\_feasible} (TUNE\_PER 1, 5, 25); random forest \texttt{all\_features} (TUNE\_PER 1, 5, 25)}
\newcommand{\campTPRuleN}{9}
\newcommand{\campTPRuleQBelow}{2}
\newcommand{\campTPRuleQAbove}{0}
\newcommand{\campTPRuleSBelow}{1}
\newcommand{\campTPRuleSAbove}{0}
\newcommand{\campOptTN}{36}
\newcommand{\campOptTQBelow}{0}
\newcommand{\campOptTQAbove}{19}
\newcommand{\campOptTSBelow}{0}
\newcommand{\campOptTSAbove}{0}
\newcommand{\campOptRSN}{36}
\newcommand{\campOptRSQBelow}{8}
\newcommand{\campOptRSQAbove}{3}
\newcommand{\campOptRSSBelow}{1}
\newcommand{\campOptRSSAbove}{0}
\newcommand{\campOptRidgeN}{36}
\newcommand{\campOptRidgeQBelow}{0}
\newcommand{\campOptRidgeQAbove}{14}
\newcommand{\campOptRidgeSBelow}{2}
\newcommand{\campOptRidgeSAbove}{0}
\newcommand{\campOptRSBelowWho}{XGBoost \texttt{live\_feasible} (TUNE\_PER 250); XGBoost \texttt{all\_features} (TUNE\_PER 250); random forest \texttt{live\_feasible} (TUNE\_PER 1, 5, 25); random forest \texttt{all\_features} (TUNE\_PER 1, 5, 25)}
\newcommand{\campPickLgbmMcsLo}{14}
\newcommand{\campPickLgbmMcsHi}{23}
\newcommand{\campPickRfDepthLo}{20}
\newcommand{\campPickRfDepthHi}{28}
\newcommand{\campPickRfDepthPrior}{17}
\newcommand{\campPickLambdaHighMax}{0.0}
\newcommand{\campPickLambdaPriorLo}{5.1}
\newcommand{\campLstmiVsLSTMN}{3}
\newcommand{\campLstmiVsLSTMQBelow}{1}
\newcommand{\campLstmiVsLSTMQAbove}{0}
\newcommand{\campLstmiVsLSTMSBelow}{0}
\newcommand{\campLstmiVsLSTMSAbove}{0}
\newcommand{\campLstmiVsRidgeN}{3}
\newcommand{\campLstmiVsRidgeQBelow}{0}
\newcommand{\campLstmiVsRidgeQAbove}{3}
\newcommand{\campLstmiVsRidgeSBelow}{0}
\newcommand{\campLstmiVsRidgeSAbove}{0}
\newcommand{\campLstmiStepBase}{19}
\newcommand{\campLstmiKeptBase}{19}
\newcommand{\campLstmiKeptBaseLo}{16}
\newcommand{\campLstmiKeptBaseHi}{19}
\newcommand{\campLstmiStepLive}{54}
\newcommand{\campLstmiKeptLive}{40}
\newcommand{\campLstmiKeptLiveLo}{37}
\newcommand{\campLstmiKeptLiveHi}{41}
\newcommand{\campLstmiStepAll}{120}
\newcommand{\campLstmiKeptAll}{80}
\newcommand{\campLstmiKeptAllLo}{74}
\newcommand{\campLstmiKeptAllHi}{82}
\newcommand{\campLstmiTunings}{18}
\newcommand{\campLstmiNThirteen}{3}
\newcommand{\campLstmiNFortyEight}{9}
\newcommand{\campLstmiNNinetySix}{6}
\newcommand{\campLstmiRefits}{1{,}469}
\newcommand{\campXColsLive}{28}
\newcommand{\campXColsAll}{94}
\newcommand{\campXMaxAbs}{\ensuremath{2.6\times10^{-11}}}
\newcommand{\campXHosts}{12}
\newcommand{\campXSamePairs}{18}
\newcommand{\campSameClassTrials}{900/900}
\newcommand{\campSameClassTrialsRep}{900/900}
\newcommand{\campSameClassFcst}{1{,}620/1{,}620}
\newcommand{\campXTrialsLgbm}{100/300}
\newcommand{\campXFcstLgbm}{5/360}
\newcommand{\campXFcstMaxLgbm}{0.103}
\newcommand{\campXPicksLgbm}{not identical}
\newcommand{\campXTrialsXgb}{211/300}
\newcommand{\campXFcstXgb}{344/360}
\newcommand{\campXFcstMaxXgb}{0.033}
\newcommand{\campXPicksXgb}{identical}
\newcommand{\campXTrialsRf}{283/300}
\newcommand{\campXFcstRf}{360/360}
\newcommand{\campXFcstMaxRf}{0.000}
\newcommand{\campXPicksRf}{identical}
\newcommand{\campXcPoints}{20}
\newcommand{\campXcPickLgbm}{18}
\newcommand{\campXcPickXgb}{0}
\newcommand{\campXcPickRf}{0}
\newcommand{\campXcValLgbm}{3.6}
\newcommand{\campXcValXgb}{\ensuremath{3.9\times10^{-3}}}
\newcommand{\campXcFcstXgb}{\ensuremath{7.3\times10^{-3}}}
\newcommand{\campXcSame}{20}
\newcommand{\campClassRelTen}{4.5}
\newcommand{\campClassQTen}{\ensuremath{-}0.1}
\newcommand{\campClassQTenCI}{[\ensuremath{-}0.8, \ensuremath{+}0.5]}
\newcommand{\campClassSTen}{\ensuremath{-}0.26}
\newcommand{\campClassSTenCI}{[\ensuremath{-}0.68, \ensuremath{+}0.11]}
\newcommand{\campClassRelOne}{5.5}
\newcommand{\campClassQOne}{\ensuremath{-}0.1}
\newcommand{\campClassQOneCI}{[\ensuremath{-}0.8, \ensuremath{+}0.8]}
\newcommand{\campClassSOne}{\ensuremath{+}0.07}
\newcommand{\campClassSOneCI}{[\ensuremath{-}0.30, \ensuremath{+}0.42]}
\newcommand{\campClassRows}{1{,}469}
\newcommand{\campLstmiXTune}{17/18}
\newcommand{\campLstmiXChunk}{19/20}
\newcommand{\campSMDiffer}{35}
\newcommand{\campSMMaxLgbm}{0.65}
\newcommand{\campSMMaxXgb}{0.17}
\newcommand{\campSMMaxRf}{0.0025}
\newcommand{\campSMN}{81}
\newcommand{\campSMQBelow}{2}
\newcommand{\campSMQAbove}{3}
\newcommand{\campSMSBelow}{0}
\newcommand{\campSMSAbove}{0}
\newcommand{\campSMSharpeExcl}{0}
\newcommand{\campCChunks}{214}
\newcommand{\campCChunkClasses}{1}
\newcommand{\campHChunks}{327}
\newcommand{\campHChunkClasses}{1}
\newcommand{\campHGateChunks}{6}
\newcommand{\campMTKeysBoth}{136}
\newcommand{\campMTKeysNew}{114}
\newcommand{\campMTHeadQ}{0.1005}
\newcommand{\campMTHeadQBefore}{0.1005}
\newcommand{\campMTHeadPos}{0}
\newcommand{\campMTHeadDays}{866}
\newcommand{\campMTHeadS}{1.90}
\newcommand{\campMTHeadSX}{1.44}
\newcommand{\campMTHeadMaxRel}{\ensuremath{1.0\times10^{-12}}}
\newcommand{\campMTSigWho}{per-bar XGBoost [live\_feasible]: 1.70 $\to$ 1.27, \ensuremath{-}0.43 [\ensuremath{-}0.94, \ensuremath{-}0.03]; per-bar XGBoost [baseline]: 1.12 $\to$ 0.92, \ensuremath{-}0.21 [\ensuremath{-}0.42, \ensuremath{-}0.04]; per-bar LSTM [baseline]: 0.61 $\to$ 1.26, \ensuremath{+}0.65 [\ensuremath{+}0.05, \ensuremath{+}1.27]}
\newcommand{\campMTSigN}{3}
\newcommand{\campMTNOther}{219}
\newcommand{\campMTBeatMid}{0}
\newcommand{\campMTBeatCrossed}{0}
\newcommand{\campMTWorseMid}{43}
\newcommand{\campMTQBetter}{3}
\newcommand{\campMTQWorse}{37}
\newcommand{\campMTQBetterFamily}{HAR-ladder variants}
\newcommand{\campMTQBetterWho}{per-bar elastic net [live\_feasible], HAR ladder base3: QLIKE 0.0934 (\ensuremath{-}7.1\,\%, interval on the daily difference [\ensuremath{-}0.0144, \ensuremath{-}0.0019]), $\Delta$Sharpe mid \ensuremath{-}0.33 [\ensuremath{-}0.91, \ensuremath{+}0.34]; per-bar elastic net [live\_feasible], HAR ladder base2: QLIKE 0.0944 (\ensuremath{-}6.1\,\%, interval on the daily difference [\ensuremath{-}0.0122, \ensuremath{-}0.0012]), $\Delta$Sharpe mid \ensuremath{-}0.15 [\ensuremath{-}0.82, \ensuremath{+}0.59]; per-bar lasso [live\_feasible], HAR ladder base2: QLIKE 0.0948 (\ensuremath{-}5.8\,\%, interval on the daily difference [\ensuremath{-}0.0126, \ensuremath{-}0.0001]), $\Delta$Sharpe mid \ensuremath{-}0.12 [\ensuremath{-}0.76, \ensuremath{+}0.55]}
\newcommand{\campFIRidgeNoise}{1.7}
\newcommand{\campFILeadN}{54}
\newcommand{\campFILead}{52}
\newcommand{\campFISpearmanTrees}{0.90}
\newcommand{\campFISpearmanTreesMin}{0.66}
\newcommand{\campDvsClaims}{23}
\newcommand{\campDvsClaimsSame}{20}
\newcommand{\campDvsLgbmQLive}{\ensuremath{+}0.0004}
\newcommand{\campDvsLgbmQLiveCI}{[\ensuremath{-}0.0012, \ensuremath{+}0.0020]}
\newcommand{\campDvsLgbmSLive}{\ensuremath{+}0.04}
\newcommand{\campDvsLgbmSLiveCI}{[\ensuremath{-}0.30, \ensuremath{+}0.39]}
\newcommand{\campDvsLgbmQAll}{\ensuremath{+}0.0002}
\newcommand{\campDvsLgbmQAllCI}{[\ensuremath{-}0.0012, \ensuremath{+}0.0016]}
\newcommand{\campDvsLgbmSAll}{\ensuremath{+}0.20}
\newcommand{\campDvsLgbmSAllCI}{[\ensuremath{-}0.27, \ensuremath{+}0.77]}
\newcommand{\campDvsLinMaxRel}{\ensuremath{3.3\times10^{-12}}}
\newcommand{\campDvsRfLive}{13.2 $\to$ 6.8}
\newcommand{\campComputeOptunaAlloc}{1422.6}
\newcommand{\campComputeOptunaPeak}{900}


% Shrink a table to the text width only when it is wider; never enlarge it.
\newsavebox{\campbox}
\newcommand{\campfit}[1]{%
  \sbox{\campbox}{#1}%
  \ifdim\wd\campbox>\textwidth
    \resizebox{\textwidth}{!}{\usebox{\campbox}}%
  \else
    \usebox{\campbox}%
  \fi}

\subsection{The 16:00 Models Re-run: Design, Mask, Cadence, Tuning}
\label{sec:app_campaign_16h}

Every per-bar model that forecasts the 16:00 bar (the 15:30--16:00 half hour, issued at
15:30, the bar the last-30-min trade is decided on) was re-run with one change at a time:
the design (\ref{sec:app_camp_design}), a per-window column mask for the trees and the LSTM
(\ref{sec:app_camp_mask}), the refit cadence (\ref{sec:app_camp_cadence}), the tuning
(\ref{sec:app_camp_tuning}); an intraday-sequence LSTM was added
(\ref{sec:app_camp_lstmi}). The section also records what the re-run showed about
reproducibility across CPUs (\ref{sec:app_camp_repro}), the master table before and after
(\ref{sec:app_camp_master}), the importance and density studies on the new design
(\ref{sec:app_camp_featimp}) and the compute used (\ref{sec:app_camp_compute}). Nothing is
discarded: a rung that did worse is reported beside one that did better.

Every comparison uses the research scorer (the 16:00 bar recalibrated on its own), the same
\campTradeDays{} trade days, and the $\mathrm{sign}(s)$ rule on the straddle (nearest
out-of-the-money call and put, same-day expiry, one position). A difference is $a-b$; a
QLIKE difference is given in percent of $b$'s QLIKE, with its 95\,\% paired day-block
bootstrap interval in the same units (negative = $a$ has the lower loss); a Sharpe
difference is the midpoint-fill $\mathrm{sign}(s)$ Sharpe ratio of $a$ minus that of $b$,
with its interval. In the grid tables each cell holds the QLIKE difference on its first line
and the Sharpe difference on its second. Input sets: \texttt{baseline} (the HAR ladder and
the calendar columns), \texttt{live\_feasible} (the inputs a 15:30 forecaster can rebuild
live) and \texttt{all\_features}. Key to the design columns named below:
\texttt{har\_ma\_*} is the bar's realized variance averaged over the last 1, 5, 25, 125,
625 and 3125 bars (the HAR ladder, \texttt{har\_ma\_k} one rung);
\texttt{har\_ma\_k\_x\_open} and \texttt{har\_ma\_k\_x\_close} are that rung times the
session-open and session-close indicators (the session-edge interactions).

% ----------------------------------------------------------------------------
\subsubsection{The design: the session-edge columns}
\label{sec:app_camp_design}

The session-edge interactions were built for the pooled 48-bar models. In a
one-bar-per-session series at the 16:00 bar the \campNSessionEdgeHalf{}
\texttt{har\_ma\_k\_x\_open} columns are identically zero and the \campNSessionEdgeHalf{}
\texttt{har\_ma\_k\_x\_close} columns are exact copies of \texttt{har\_ma\_k}. The
\campNSessionEdge{} columns were dropped from the per-bar design and kept in the pooled
models: \texttt{baseline} goes from \campPOldBase{} to \campPNewBase{} columns,
\texttt{live\_feasible} from \campPOldLive{} to \campPNewLive{}, \texttt{all\_features} from
\campPOldAll{} to \campPNewAll{} (Table~\ref{tab:app_camp_design}).

The linear arms' identifiability mask --- at every penalty tune, drop the columns constant
on the tuning window and the exact copies of an earlier kept column --- had already removed
all of them: in \campMaskDropAgree{} of the \campNLinearMasked{} masked linear forecasts the
number of masked columns per tune falls by exactly the number of dropped columns, and the
columns each fit keeps are the same before and after (last column of
Table~\ref{tab:app_camp_design}). The fits solve the same problem, and the \campNLinear{}
linear forecasts of the master table move by rounding only
(Table~\ref{tab:app_camp_linear}): OLS by at most \campOLSMaxRel{} and ridge by at most
\campRidgeMaxRel{} (relative change of the 16:00 forecast). The lasso and the elastic net
move by up to \campLassoMaxRel{} and \campEnetMaxRel{}, on the rows of one penalty-tune
period, in \campLassoMoved{} of \campLassoN{} and \campEnetMoved{} of \campEnetN{} forecasts;
positions change on \campPosChanged{} of the \campNLinear{}\,$\times$\,\campTradeDays{}
forecast-days, and no Sharpe ratio moves by more than \campMaxDSharpeLin{}. A control on the
\campAttrN{} arms that moved attributes this to the solvers' floating-point path: the
pre-change code re-run on the same CPU architecture reproduces its first run bit for bit in
\campAttrRepeat{} of \campAttrN{}; the new code against the pre-change code on one
architecture moves \campAttrDesignN{} of them, by \campAttrDesignLo{} to
\campAttrDesignHi{}; the pre-change code alone, moved to another machine, moves
\campAttrMachineN{}, by \campAttrMachineLo{} to \campAttrMachineHi{}. The LSTM, which had no
mask at that step, moved more (median relative change of the 16:00 forecast
\campLSTMDedupLo--\campLSTMDedupHi\,\% over the three input sets; a new input width also
redraws its initial weights).

The rest-of-day model's \campRodArms{} arms were re-run on the same design; their mask fell
by exactly \campNSessionEdge{} columns per tune in \campRodMaskSame{} of them. The master
table's \campRodChecks{} check rows (the rest-of-day model at 15:30, which should equal its
per-bar 16:00 twin) had differed from their twins by up to \campRodBeforeMax{} (the same
position on as few as \campRodBeforeSameMin{} of \campRodDeck{} days); they are now
bit-identical to their twins in \campRodAfterBit{} of \campRodChecks{}, on every one of the
\campRodRows{} 16:00 rows.

% Source: results/linear_subsection_dedup/change_by_arm.csv
\begin{table}[H]
\centering
\caption{The per-bar design before and after the session-edge columns were dropped, and
the linear arms' identifiability mask per penalty tune (range over tunes) for the per-bar
ridge: the columns it masked fall by exactly the columns removed, so the columns kept are
unchanged.}
\label{tab:app_camp_design}
\campfit{\small% AUTO-GENERATED by writeup/make_appendix_campaign_16h_tex.py -- do not edit.
% Source: results/linear_subsection_dedup/change_by_arm.csv (rows sub_ridge_<input set>)
\begin{tabular}{lrrrrrr}
\toprule
input set & \shortstack{columns\\before} & \shortstack{columns\\after} & removed & \shortstack{masked per\\tune, before} & \shortstack{masked per\\tune, after} & \shortstack{kept per\\tune} \\
\midrule
\texttt{baseline} & 34 & 22 & 12 & 17 & 5 & 17 \\
\texttt{live\_feasible} & 244 & 232 & 12 & 106--108 & 94--96 & 136--138 \\
\texttt{all\_features} & 640 & 628 & 12 & 262--280 & 250--268 & 360--378 \\
\bottomrule
\end{tabular}
}
\end{table}

% Source: results/linear_subsection_dedup/change_by_arm.csv, gates.csv
\begin{table}[H]
\centering
\caption{The master table's per-bar linear forecasts, new design minus old, by estimator:
the largest and the median (over forecasts) of the per-forecast maximum and median relative
change of the 16:00 forecast, the forecasts whose maximum exceeds the reproduction gate's
bound, the trade positions changed (forecast-days, of \campTradeDays{} per forecast) and the
largest change of a $\mathrm{sign}(s)$ Sharpe ratio (midpoint).}
\label{tab:app_camp_linear}
\campfit{\small% AUTO-GENERATED by writeup/make_appendix_campaign_16h_tex.py -- do not edit.
% Source: results/linear_subsection_dedup/change_by_arm.csv
% Source: results/linear_subsection_dedup/gates.csv (the 'moved' bound)
\begin{tabular}{lrrrrrr}
\toprule
estimator & forecasts & \shortstack{max rel.\\change} & \shortstack{median rel.\\change} & \shortstack{moved\\(> \ensuremath{1\times10^{-8}})} & \shortstack{positions\\changed} & \shortstack{max $|\Delta$Sharpe$|$\\mid} \\
\midrule
OLS & 1 & \ensuremath{8.4\times10^{-15}} & \ensuremath{4.4\times10^{-16}} & 0 & 0 & 0.00 \\
ridge & 25 & \ensuremath{4.3\times10^{-10}} & \ensuremath{1.1\times10^{-14}} & 0 & 0 & 0.00 \\
lasso & 25 & \ensuremath{9.1\times10^{-2}} & \ensuremath{2.8\times10^{-15}} & 7 & 6 & 0.06 \\
elastic net & 25 & \ensuremath{9.5\times10^{-2}} & \ensuremath{3.0\times10^{-15}} & 10 & 3 & 0.04 \\
\bottomrule
\end{tabular}
}
\end{table}

% ----------------------------------------------------------------------------
\subsubsection{One mask for every model}
\label{sec:app_camp_mask}

The trees and the LSTM were then given the linear arms' rule. At every fit --- a refit on
its training window, or a tuning point on the window it tunes on --- every column constant on
that window and every exact copy of an earlier kept column is dropped, first occurrences
kept in column order; the LSTM, whose network is rebuilt at every refit, recomputes it at
every refit. Importance and TreeSHAP are mapped back to all columns with zero for a dropped
column. The rule looks at the window only, so every model of an input set keeps the same
columns at the same refit: over the \campKeptRefits{} daily refits \texttt{baseline} keeps
\campKeptBase{} of \campPBase{} columns at every refit, \texttt{live\_feasible}
\campKeptLive{} of \campPLive{} (median; range \campKeptLiveLo--\campKeptLiveHi), and
\texttt{all\_features} \campKeptAll{} of \campPAll{} (\campKeptAllLo--\campKeptAllHi); the
Optuna runs below keep the same counts.

Table~\ref{tab:app_camp_mask} sets each masked rung against the same rung unmasked (both
on the new design). Rungs: T10 and T1 are the shipped tree configurations refit every 10
sessions and every session; RS10 and RS1 the random search, refit every 10 sessions and
every session; LSTM10 the per-bar LSTM refit every 10 sessions. Over the \campMaskTreeN{}
tree rows the median absolute QLIKE change is \campMaskMedAbsQ\,\%; \campMaskTreeQAbove{}
QLIKE intervals lie above zero (the masked rung has the higher loss) and
\campMaskTreeQBelow{} below; \campMaskTreeSBelow{} Sharpe interval lies below zero and
\campMaskTreeSAbove{} above. For the LSTM \campMaskLSTMQBelow{} of \campMaskLSTMN{} QLIKE
intervals lie below zero (\texttt{all\_features}, \campMaskLSTMAllQ\,\%
\campMaskLSTMAllQCI) and \campMaskLSTMQAbove{} above. The masked rungs ran on one CPU class;
the unmasked rungs did not pin the class (the class of their chunks was not recorded), so a
masked-minus-unmasked difference also contains any difference of CPU class, whose size alone
is recorded in \ref{sec:app_camp_repro}.

% Source: results/trees_mask_1600/mask_pairs.csv (family 'masked - unmasked')
\begin{table}[H]
\centering
\caption{The per-window mask, masked minus unmasked, per rung, model and input set, same
design, \campTradeDays{} days. Each cell: QLIKE difference in \% of the unmasked QLIKE
[95\,\% interval] over the $\mathrm{sign}(s)$ Sharpe difference at the midpoint [95\,\%
interval]. The unmasked rungs were not pinned to a CPU class.}
\label{tab:app_camp_mask}
\campfit{\scriptsize% AUTO-GENERATED by writeup/make_appendix_campaign_16h_tex.py -- do not edit.
% Source: results/trees_mask_1600/mask_pairs.csv (family 'masked - unmasked')
\begin{tabular}{llccccc}
\toprule
 & model & \shortstack{T10\\masked $-$ unmasked} & \shortstack{T1\\masked $-$ unmasked} & \shortstack{RS10\\masked $-$ unmasked} & \shortstack{RS1\\masked $-$ unmasked} & \shortstack{LSTM10\\masked $-$ unmasked} \\
\midrule
\multicolumn{7}{l}{\emph{\texttt{baseline}}} \\
 & LightGBM & \shortstack{\ensuremath{+}0.5 [\ensuremath{-}0.6, \ensuremath{+}1.5]\\\ensuremath{-}0.18 [\ensuremath{-}0.80, \ensuremath{+}0.30]} & \shortstack{0.0 [\ensuremath{-}0.9, \ensuremath{+}0.9]\\\ensuremath{-}0.25 [\ensuremath{-}0.74, \ensuremath{+}0.15]} & \shortstack{\ensuremath{-}0.2 [\ensuremath{-}1.1, \ensuremath{+}0.6]\\\ensuremath{+}0.08 [\ensuremath{-}0.09, \ensuremath{+}0.27]} & \shortstack{\ensuremath{-}0.2 [\ensuremath{-}1.1, \ensuremath{+}0.6]\\\ensuremath{-}0.04 [\ensuremath{-}0.13, \ensuremath{+}0.04]} &  \\
 & XGBoost & \shortstack{\ensuremath{+}0.4 [\ensuremath{-}0.3, \ensuremath{+}1.0]\\\ensuremath{-}0.01 [\ensuremath{-}0.21, \ensuremath{+}0.16]} & \shortstack{\ensuremath{+}0.3 [\ensuremath{-}0.3, \ensuremath{+}0.9]\\\ensuremath{-}0.07 [\ensuremath{-}0.30, \ensuremath{+}0.11]} & \shortstack{\ensuremath{-}0.2 [\ensuremath{-}1.9, \ensuremath{+}1.4]\\\ensuremath{-}0.03 [\ensuremath{-}0.48, \ensuremath{+}0.41]} & \shortstack{\ensuremath{-}0.7 [\ensuremath{-}2.3, \ensuremath{+}1.0]\\\ensuremath{+}0.16 [\ensuremath{-}0.06, \ensuremath{+}0.40]} &  \\
 & random forest & \shortstack{\ensuremath{+}0.4 [\ensuremath{-}0.3, \ensuremath{+}1.2]\\0.00 [\ensuremath{-}0.25, \ensuremath{+}0.25]} & \shortstack{0.0 [\ensuremath{-}0.7, \ensuremath{+}0.7]\\\ensuremath{+}0.06 [\ensuremath{-}0.16, \ensuremath{+}0.31]} & \shortstack{\ensuremath{-}0.1 [\ensuremath{-}1.5, \ensuremath{+}1.4]\\\ensuremath{+}0.16 [\ensuremath{-}0.46, \ensuremath{+}0.79]} & \shortstack{\ensuremath{+}0.4 [\ensuremath{-}1.1, \ensuremath{+}1.9]\\\ensuremath{+}0.32 [\ensuremath{-}0.26, \ensuremath{+}0.88]} &  \\
 & LSTM &  &  &  &  & \shortstack{\ensuremath{+}0.2 [\ensuremath{-}2.9, \ensuremath{+}4.2]\\\ensuremath{-}0.03 [\ensuremath{-}0.55, \ensuremath{+}0.49]} \\
\multicolumn{7}{l}{\emph{\texttt{live\_feasible}}} \\
 & LightGBM & \shortstack{\ensuremath{+}1.0 [\ensuremath{-}0.4, \ensuremath{+}2.3]\\\ensuremath{-}0.17 [\ensuremath{-}0.80, \ensuremath{+}0.39]} & \shortstack{\ensuremath{-}1.0 [\ensuremath{-}2.4, \ensuremath{+}0.3]\\\ensuremath{-}0.06 [\ensuremath{-}0.49, \ensuremath{+}0.37]} & \shortstack{\ensuremath{-}0.9 [\ensuremath{-}2.7, \ensuremath{+}0.8]\\\ensuremath{-}0.12 [\ensuremath{-}0.67, \ensuremath{+}0.43]} & \shortstack{\ensuremath{-}0.1 [\ensuremath{-}2.1, \ensuremath{+}1.8]\\\ensuremath{+}0.01 [\ensuremath{-}0.55, \ensuremath{+}0.58]} &  \\
 & XGBoost & \shortstack{\ensuremath{+}0.3 [\ensuremath{-}0.5, \ensuremath{+}1.3]\\\ensuremath{-}0.05 [\ensuremath{-}0.37, \ensuremath{+}0.33]} & \shortstack{\ensuremath{+}1.3 [\ensuremath{+}0.5, \ensuremath{+}2.3]\\\ensuremath{-}0.36 [\ensuremath{-}0.73, \ensuremath{-}0.06]} & \shortstack{\ensuremath{-}1.1 [\ensuremath{-}3.6, \ensuremath{+}1.2]\\\ensuremath{+}0.16 [\ensuremath{-}0.21, \ensuremath{+}0.60]} & \shortstack{\ensuremath{-}1.1 [\ensuremath{-}4.0, \ensuremath{+}1.4]\\\ensuremath{-}0.01 [\ensuremath{-}0.29, \ensuremath{+}0.28]} &  \\
 & random forest & \shortstack{\ensuremath{+}0.2 [\ensuremath{-}0.8, \ensuremath{+}1.1]\\\ensuremath{-}0.08 [\ensuremath{-}0.36, \ensuremath{+}0.20]} & \shortstack{\ensuremath{+}0.1 [\ensuremath{-}0.9, \ensuremath{+}1.1]\\\ensuremath{+}0.10 [\ensuremath{-}0.17, \ensuremath{+}0.44]} & \shortstack{\ensuremath{+}12.8 [\ensuremath{-}0.7, \ensuremath{+}32.8]\\\ensuremath{+}0.37 [\ensuremath{-}0.48, \ensuremath{+}1.24]} & \shortstack{\ensuremath{+}14.7 [\ensuremath{-}0.2, \ensuremath{+}37.1]\\\ensuremath{-}0.09 [\ensuremath{-}0.77, \ensuremath{+}0.71]} &  \\
 & LSTM &  &  &  &  & \shortstack{\ensuremath{+}0.3 [\ensuremath{-}3.0, \ensuremath{+}3.4]\\\ensuremath{-}0.05 [\ensuremath{-}0.69, \ensuremath{+}0.57]} \\
\multicolumn{7}{l}{\emph{\texttt{all\_features}}} \\
 & LightGBM & \shortstack{\ensuremath{-}0.7 [\ensuremath{-}2.3, \ensuremath{+}1.1]\\\ensuremath{+}0.09 [\ensuremath{-}0.34, \ensuremath{+}0.53]} & \shortstack{\ensuremath{+}1.5 [\ensuremath{-}0.3, \ensuremath{+}3.4]\\\ensuremath{-}0.16 [\ensuremath{-}0.73, \ensuremath{+}0.35]} & \shortstack{\ensuremath{+}2.7 [\ensuremath{-}0.4, \ensuremath{+}6.2]\\\ensuremath{+}0.06 [\ensuremath{-}0.59, \ensuremath{+}0.63]} & \shortstack{\ensuremath{+}0.7 [\ensuremath{-}3.5, \ensuremath{+}4.2]\\\ensuremath{-}0.30 [\ensuremath{-}0.95, \ensuremath{+}0.34]} &  \\
 & XGBoost & \shortstack{\ensuremath{+}0.9 [\ensuremath{-}0.2, \ensuremath{+}2.1]\\\ensuremath{-}0.36 [\ensuremath{-}0.87, \ensuremath{+}0.08]} & \shortstack{\ensuremath{+}0.1 [\ensuremath{-}0.7, \ensuremath{+}1.0]\\\ensuremath{+}0.33 [\ensuremath{-}0.11, \ensuremath{+}0.80]} & \shortstack{\ensuremath{+}6.5 [\ensuremath{+}1.8, \ensuremath{+}11.6]\\\ensuremath{-}0.44 [\ensuremath{-}0.92, \ensuremath{+}0.04]} & \shortstack{\ensuremath{+}6.2 [\ensuremath{+}1.6, \ensuremath{+}11.3]\\\ensuremath{-}0.36 [\ensuremath{-}0.85, \ensuremath{+}0.09]} &  \\
 & random forest & \shortstack{\ensuremath{+}0.1 [\ensuremath{-}0.8, \ensuremath{+}1.0]\\\ensuremath{-}0.27 [\ensuremath{-}0.62, \ensuremath{+}0.06]} & \shortstack{\ensuremath{+}0.9 [0.0, \ensuremath{+}1.7]\\\ensuremath{-}0.17 [\ensuremath{-}0.70, \ensuremath{+}0.29]} & \shortstack{\ensuremath{+}9.4 [\ensuremath{+}0.9, \ensuremath{+}23.1]\\\ensuremath{+}0.31 [\ensuremath{-}0.29, \ensuremath{+}0.94]} & \shortstack{\ensuremath{+}10.1 [\ensuremath{+}0.6, \ensuremath{+}26.5]\\\ensuremath{+}0.16 [\ensuremath{-}0.29, \ensuremath{+}0.72]} &  \\
 & LSTM &  &  &  &  & \shortstack{\ensuremath{-}17.1 [\ensuremath{-}33.2, \ensuremath{-}4.2]\\\ensuremath{-}0.01 [\ensuremath{-}0.84, \ensuremath{+}0.83]} \\
\bottomrule
\end{tabular}
}
\end{table}

% ----------------------------------------------------------------------------
\subsubsection{Refit cadence}
\label{sec:app_camp_cadence}

With the mask in place, the trees refit every session, as the per-bar linear arms do, and
every 10th session as before (first two columns of Table~\ref{tab:app_camp_cadence}). For
the shipped configuration, T1 minus T10 ranges from \campTOneTenQLo\,\% to
\campTOneTenQHi\,\% in QLIKE, with \campTOneTenQBelow{} of \campTOneTenN{} intervals below
zero and \campTOneTenQAbove{} above; \campTOneTenSAbove{} Sharpe interval lies above zero
and \campTOneTenSBelow{} below. For the random search, RS1 minus RS10 ranges from
\campRSOneTenQLo\,\% to \campRSOneTenQHi\,\%, with \campRSOneTenQBelow{} of
\campRSOneTenN{} QLIKE intervals below zero and \campRSOneTenQAbove{} above; no Sharpe
interval excludes zero (\campRSOneTenSBelow{} below, \campRSOneTenSAbove{} above). The
per-bar LSTM refit every session against every 10th: \campLSTMCadBase\,\%
\campLSTMCadBaseCI{} (\texttt{baseline}), \campLSTMCadLive\,\% \campLSTMCadLiveCI{}
(\texttt{live\_feasible}), \campLSTMCadAll\,\% \campLSTMCadAllCI{} (\texttt{all\_features});
\campLSTMCadQBelow{} of \campLSTMCadN{} intervals below zero, and no Sharpe interval
excludes zero (\campLSTMCadSBelow{} below, \campLSTMCadSAbove{} above).

Against the per-bar ridge, none of the \campRidgeTreeN{} masked tree rows has a QLIKE
interval below zero (\campRidgeTreeQBelow{}) and \campRidgeTreeQAbove{} have one above; the
lowest relative QLIKE is \campRidgeBestTree{} at \campRidgeBestTreeQ\,\%
\campRidgeBestTreeQCI; no Sharpe interval excludes zero (\campRidgeTreeSBelow{} below,
\campRidgeTreeSAbove{} above). The daily-refit LSTM minus the ridge is
\campRidgeLSTMBase\,\% in QLIKE on \texttt{baseline}, \campRidgeLSTMLive\,\% on
\texttt{live\_feasible} and \campRidgeLSTMAll\,\% on \texttt{all\_features}, with
\campRidgeLSTMQAbove{} of \campRidgeLSTMN{} intervals above zero.

% Source: results/trees_mask_1600/mask_pairs.csv (family 'masked ladder')
\begin{table}[H]
\centering
\caption{The masked ladder: refit every session minus every 10th (shipped configuration;
random search; the per-bar LSTM in the first column), and random search minus the shipped
configuration, both refit every session. Cells as in Table~\ref{tab:app_camp_mask}, the
second forecast of each pair being the comparator.}
\label{tab:app_camp_cadence}
\campfit{\scriptsize% AUTO-GENERATED by writeup/make_appendix_campaign_16h_tex.py -- do not edit.
% Source: results/trees_mask_1600/mask_pairs.csv (family 'masked ladder')
\begin{tabular}{llccc}
\toprule
 & model & \shortstack{T1 $-$ T10\\(refit 1 vs 10)} & \shortstack{RS1 $-$ RS10\\(refit 1 vs 10)} & \shortstack{RS1 $-$ T1\\(search vs shipped)} \\
\midrule
\multicolumn{5}{l}{\emph{\texttt{baseline}}} \\
 & LightGBM & \shortstack{\ensuremath{-}0.8 [\ensuremath{-}2.3, \ensuremath{+}0.7]\\\ensuremath{+}0.18 [\ensuremath{-}0.19, \ensuremath{+}0.52]} & \shortstack{\ensuremath{+}0.1 [\ensuremath{-}0.8, \ensuremath{+}1.0]\\\ensuremath{-}0.29 [\ensuremath{-}0.76, \ensuremath{+}0.09]} & \shortstack{\ensuremath{+}7.2 [\ensuremath{+}2.1, \ensuremath{+}14.0]\\\ensuremath{-}0.28 [\ensuremath{-}0.89, \ensuremath{+}0.30]} \\
 & XGBoost & \shortstack{\ensuremath{-}0.8 [\ensuremath{-}1.6, 0.0]\\\ensuremath{+}0.03 [\ensuremath{-}0.16, \ensuremath{+}0.29]} & \shortstack{\ensuremath{-}0.8 [\ensuremath{-}1.7, \ensuremath{+}0.2]\\\ensuremath{+}0.10 [\ensuremath{-}0.11, \ensuremath{+}0.38]} & \shortstack{\ensuremath{+}5.2 [\ensuremath{+}0.2, \ensuremath{+}12.8]\\\ensuremath{+}0.27 [\ensuremath{-}0.18, \ensuremath{+}0.76]} \\
 & random forest & \shortstack{\ensuremath{-}1.3 [\ensuremath{-}3.3, \ensuremath{+}0.8]\\\ensuremath{+}0.27 [\ensuremath{-}0.42, \ensuremath{+}1.07]} & \shortstack{\ensuremath{-}0.7 [\ensuremath{-}1.7, \ensuremath{+}0.3]\\\ensuremath{+}0.19 [\ensuremath{-}0.20, \ensuremath{+}0.57]} & \shortstack{\ensuremath{+}4.6 [\ensuremath{-}2.2, \ensuremath{+}14.1]\\\ensuremath{+}0.08 [\ensuremath{-}0.55, \ensuremath{+}0.67]} \\
 & LSTM (1 vs 10) & \shortstack{\ensuremath{-}1.9 [\ensuremath{-}3.5, \ensuremath{-}0.3]\\\ensuremath{+}0.09 [\ensuremath{-}0.17, \ensuremath{+}0.34]} &  &  \\
\multicolumn{5}{l}{\emph{\texttt{live\_feasible}}} \\
 & LightGBM & \shortstack{\ensuremath{-}3.6 [\ensuremath{-}7.4, \ensuremath{-}0.8]\\\ensuremath{-}0.17 [\ensuremath{-}1.11, \ensuremath{+}0.59]} & \shortstack{\ensuremath{+}0.2 [\ensuremath{-}1.5, \ensuremath{+}1.6]\\\ensuremath{+}0.16 [\ensuremath{-}0.19, \ensuremath{+}0.58]} & \shortstack{\ensuremath{+}4.1 [\ensuremath{-}0.2, \ensuremath{+}8.9]\\\ensuremath{+}0.21 [\ensuremath{-}0.61, \ensuremath{+}1.07]} \\
 & XGBoost & \shortstack{\ensuremath{-}0.9 [\ensuremath{-}3.0, \ensuremath{+}0.8]\\\ensuremath{+}0.26 [\ensuremath{-}0.05, \ensuremath{+}0.61]} & \shortstack{\ensuremath{-}1.9 [\ensuremath{-}5.2, \ensuremath{+}0.1]\\\ensuremath{-}0.13 [\ensuremath{-}0.52, \ensuremath{+}0.29]} & \shortstack{\ensuremath{+}6.5 [\ensuremath{+}1.2, \ensuremath{+}13.0]\\\ensuremath{-}0.19 [\ensuremath{-}0.77, \ensuremath{+}0.36]} \\
 & random forest & \shortstack{\ensuremath{-}3.0 [\ensuremath{-}6.5, \ensuremath{+}0.2]\\\ensuremath{-}0.19 [\ensuremath{-}0.88, \ensuremath{+}0.48]} & \shortstack{\ensuremath{-}2.2 [\ensuremath{-}4.9, \ensuremath{+}0.3]\\\ensuremath{+}0.16 [\ensuremath{-}0.32, \ensuremath{+}0.62]} & \shortstack{\ensuremath{+}17.9 [\ensuremath{+}2.2, \ensuremath{+}41.5]\\\ensuremath{+}0.09 [\ensuremath{-}0.66, \ensuremath{+}0.91]} \\
 & LSTM (1 vs 10) & \shortstack{\ensuremath{-}2.5 [\ensuremath{-}7.8, \ensuremath{+}3.2]\\\ensuremath{-}0.17 [\ensuremath{-}0.80, \ensuremath{+}0.48]} &  &  \\
\multicolumn{5}{l}{\emph{\texttt{all\_features}}} \\
 & LightGBM & \shortstack{\ensuremath{-}2.5 [\ensuremath{-}7.2, \ensuremath{+}0.9]\\\ensuremath{-}0.18 [\ensuremath{-}1.14, \ensuremath{+}0.77]} & \shortstack{\ensuremath{-}1.7 [\ensuremath{-}5.1, \ensuremath{+}0.7]\\\ensuremath{-}0.27 [\ensuremath{-}0.61, \ensuremath{+}0.06]} & \shortstack{\ensuremath{+}3.6 [\ensuremath{-}0.4, \ensuremath{+}8.5]\\\ensuremath{-}0.26 [\ensuremath{-}1.05, \ensuremath{+}0.61]} \\
 & XGBoost & \shortstack{\ensuremath{-}2.9 [\ensuremath{-}5.7, \ensuremath{-}1.0]\\\ensuremath{+}0.52 [\ensuremath{+}0.10, \ensuremath{+}1.02]} & \shortstack{\ensuremath{-}1.8 [\ensuremath{-}3.6, \ensuremath{-}0.4]\\\ensuremath{+}0.04 [\ensuremath{-}0.34, \ensuremath{+}0.45]} & \shortstack{\ensuremath{+}10.7 [\ensuremath{+}4.6, \ensuremath{+}17.5]\\\ensuremath{-}0.47 [\ensuremath{-}1.06, \ensuremath{+}0.11]} \\
 & random forest & \shortstack{\ensuremath{-}2.0 [\ensuremath{-}5.5, \ensuremath{+}1.1]\\\ensuremath{+}0.23 [\ensuremath{-}0.29, \ensuremath{+}0.73]} & \shortstack{\ensuremath{-}2.1 [\ensuremath{-}5.7, \ensuremath{+}0.3]\\\ensuremath{-}0.14 [\ensuremath{-}0.61, \ensuremath{+}0.28]} & \shortstack{\ensuremath{+}12.6 [\ensuremath{+}1.4, \ensuremath{+}30.0]\\\ensuremath{+}0.17 [\ensuremath{-}0.59, \ensuremath{+}0.98]} \\
 & LSTM (1 vs 10) & \shortstack{\ensuremath{-}10.1 [\ensuremath{-}18.0, \ensuremath{-}2.0]\\\ensuremath{+}0.37 [\ensuremath{-}0.39, \ensuremath{+}1.22]} &  &  \\
\bottomrule
\end{tabular}
}
\end{table}

% ----------------------------------------------------------------------------
\subsubsection{Tuning: random search and Optuna}
\label{sec:app_camp_tuning}

\paragraph{Random search against the shipped configuration.} The random search re-chooses
each tree's configuration from a fixed set of candidates on the validation tail of its
window. Both refit every session, RS1 minus T1 ranges from \campRSVsTQLo\,\% to
\campRSVsTQHi\,\% in QLIKE, with the interval above zero in \campRSVsTQAbove{} of
\campRSVsTN{} rows and below zero in \campRSVsTQBelow{};
no Sharpe interval excludes zero (\campRSVsTSBelow{} below, \campRSVsTSAbove{} above; last
column of Table~\ref{tab:app_camp_cadence}).

\paragraph{Optuna.} The random search was replaced by Optuna's TPE sampler: at every tuning
point a fresh study of \campOptTrials{} sequential trials (trial 1 the shipped
configuration, trials 1--10 TPE's random start-up) minimises validation MSE on the tail of
the tuning window after an embargo; the trees refit every session with the configuration in
force; a new tuning point is drawn every TUNE\_PER sessions (TUNE\_PER 1, 5, 25 or 250),
and at TUNE\_PER 1 and 25 the best of the first 10 and the first 25 trials is kept beside
the best of all \campOptTrials{}. Every Optuna run is masked and ran on one CPU class.

\paragraph{Are \campOptTrials{} trials enough?} Over the \campOptPoints{} tuning points of
each arm (Table~\ref{tab:app_camp_budget}), the best of the \campOptTrials{} trials falls in
trials 41--50 at \campOptLateLo--\campOptLateHi\,\% of the points, against
\campOptExch\,\% if the trials were exchangeable draws, and the shipped configuration is the
best at no more than \campOptShippedHi\,\% of them. The last 25 trials lower the best
validation MSE by \campOptGainTwentyFiveLo--\campOptGainTwentyFiveHi\,\% on average (medians
\campOptGainTwentyFiveMedLo--\campOptGainTwentyFiveMedHi\,\%) and the last 10 by
\campOptGainFortyLo--\campOptGainFortyHi\,\% (medians at most \campOptGainFortyMedMax\,\%);
the pick of the first 25 differs from the pick of all \campOptTrials{} at
\campOptPickDiffLo--\campOptPickDiffHi\,\% of the points. Out of sample
(Table~\ref{tab:app_camp_budget_oos}), of the \campBudN{} budget comparisons (best of 50
against 10, against 25, and 25 against 10, at TUNE\_PER 1 and 25) \campBudQBelow{} have a
QLIKE interval below zero and \campBudQAbove{} above; \campBudSBelow{} Sharpe interval lies
below zero and \campBudSAbove{} above.

\paragraph{How often to tune.} Against a new tuning point every 250 sessions, tuning every 1,
5 or 25 sessions (every path refit every session; Table~\ref{tab:app_camp_tuneper}) has
\campTPQBelow{} of \campTPN{} QLIKE intervals below zero --- \campTPBelowWho{} --- and
\campTPQAbove{} above; no Sharpe interval excludes zero (\campTPSBelow{} below,
\campTPSAbove{} above). Choosing the configuration by validation QLIKE instead of MSE (at
TUNE\_PER 25) has \campTPRuleQBelow{} of \campTPRuleN{} QLIKE intervals below zero and
\campTPRuleQAbove{} above, \campTPRuleSBelow{} Sharpe interval below zero and
\campTPRuleSAbove{} above.

\paragraph{Optuna against the other rungs.} Table~\ref{tab:app_camp_optuna_vs} counts the
intervals. Against the shipped configuration refit every session (T1), \campOptTQAbove{} of
\campOptTN{} Optuna QLIKE intervals lie above zero and \campOptTQBelow{} below. Against the
random search refit every session (RS1), \campOptRSQBelow{} lie below zero ---
\campOptRSBelowWho{} --- and \campOptRSQAbove{} above. Against the per-bar ridge,
\campOptRidgeQBelow{} lie below zero and \campOptRidgeQAbove{} above, with
\campOptRidgeSBelow{} Sharpe intervals below zero and \campOptRidgeSAbove{} above.

\paragraph{Where the picks sit.} Table~\ref{tab:app_camp_picks} gives, per axis, the share of
the best-of-50 picks in the outermost twentieth of the search space at either end (on the
log scale for a log axis; for the forest's depth, the smallest and the largest choice),
beside the share of TPE's random start-up draws there (the prior). The penalties
\texttt{lambda\_l2} and \texttt{reg\_lambda} are picked at their upper end in no more than
\campPickLambdaHighMax\,\% of points against a prior of at least \campPickLambdaPriorLo\,\%;
LightGBM's \texttt{min\_child\_samples} sits at its lower end in
\campPickLgbmMcsLo--\campPickLgbmMcsHi\,\% of points, and the forest picks unbounded depth in
\campPickRfDepthLo--\campPickRfDepthHi\,\% (prior \campPickRfDepthPrior\,\%).

% Source: results/linear_subsection_trees_optuna_mask/report/budget_points.csv
\begin{table}[H]
\centering
\caption{Is the Optuna budget enough? Per arm, over all tuning points: the median trial that
holds the best of 50, the share of points where it is the shipped configuration and where it
falls in trials 41--50 (\campOptExch\,\% if exchangeable), the relative gain in the best
validation MSE from the first 10, 25 and 40 trials to all 50 (mean / median), and the share
of points whose best of 25 differs from the best of 50.}
\label{tab:app_camp_budget}
\campfit{\scriptsize% AUTO-GENERATED by writeup/make_appendix_campaign_16h_tex.py -- do not edit.
% Source: results/linear_subsection_trees_optuna_mask/report/budget_points.csv
\begin{tabular}{lrrrcccr}
\toprule
arm & \shortstack{best trial\\(median)} & \shortstack{best =\\shipped (\%)} & \shortstack{best in\\41--50 (\%)} & \shortstack{val.~MSE gain\\10$\to$50 (\%)} & \shortstack{val.~MSE gain\\25$\to$50 (\%)} & \shortstack{val.~MSE gain\\40$\to$50 (\%)} & \shortstack{pick of 25\\$\neq$ of 50 (\%)} \\
\midrule
LightGBM \texttt{baseline} & 35 & 0.1 & 36.8 & 5.22 / 4.42 & 1.66 / 0.99 & 0.50 / 0.00 & 72 \\
LightGBM \texttt{live\_feasible} & 34 & 0.3 & 34.5 & 5.94 / 4.93 & 2.04 / 1.22 & 0.63 / 0.00 & 70 \\
LightGBM \texttt{all\_features} & 33 & 0.7 & 31.9 & 5.96 / 4.86 & 2.15 / 1.13 & 0.69 / 0.00 & 66 \\
XGBoost \texttt{baseline} & 36 & 0.2 & 40.2 & 4.92 / 3.94 & 1.82 / 1.13 & 0.56 / 0.00 & 74 \\
XGBoost \texttt{live\_feasible} & 34 & 0.2 & 34.7 & 5.54 / 4.70 & 1.94 / 1.14 & 0.62 / 0.00 & 71 \\
XGBoost \texttt{all\_features} & 33 & 1.0 & 33.1 & 5.24 / 4.29 & 1.95 / 1.10 & 0.63 / 0.00 & 68 \\
random forest \texttt{baseline} & 32 & 1.1 & 26.8 & 4.76 / 3.41 & 1.06 / 0.41 & 0.29 / 0.00 & 61 \\
random forest \texttt{live\_feasible} & 32 & 0.4 & 30.7 & 4.64 / 3.76 & 1.42 / 0.82 & 0.40 / 0.00 & 67 \\
random forest \texttt{all\_features} & 33 & 0.3 & 31.8 & 5.05 / 4.16 & 1.42 / 0.82 & 0.43 / 0.00 & 67 \\
\bottomrule
\end{tabular}
}
\end{table}

% Source: results/linear_subsection_trees_optuna_mask/report/budget_oos.csv
\begin{table}[H]
\centering
\caption{The budget out of sample: best of 50 trials minus best of 10 and of 25, at
TUNE\_PER 1 and 25. Cells as in Table~\ref{tab:app_camp_mask}.}
\label{tab:app_camp_budget_oos}
\campfit{\scriptsize% AUTO-GENERATED by writeup/make_appendix_campaign_16h_tex.py -- do not edit.
% Source: results/linear_subsection_trees_optuna_mask/report/budget_oos.csv
\begin{tabular}{lcccc}
\toprule
arm & \shortstack{TUNE\_PER 1:\\50 vs 10} & \shortstack{TUNE\_PER 1:\\50 vs 25} & \shortstack{TUNE\_PER 25:\\50 vs 10} & \shortstack{TUNE\_PER 25:\\50 vs 25} \\
\midrule
LightGBM \texttt{baseline} & \shortstack{\ensuremath{-}0.3 [\ensuremath{-}2.7, \ensuremath{+}2.2]\\\ensuremath{+}0.04 [\ensuremath{-}0.45, \ensuremath{+}0.48]} & \shortstack{0.0 [\ensuremath{-}1.9, \ensuremath{+}2.0]\\\ensuremath{+}0.09 [\ensuremath{-}0.19, \ensuremath{+}0.37]} & \shortstack{\ensuremath{-}0.6 [\ensuremath{-}3.0, \ensuremath{+}1.8]\\\ensuremath{+}0.17 [\ensuremath{-}0.24, \ensuremath{+}0.59]} & \shortstack{\ensuremath{+}1.4 [\ensuremath{-}0.1, \ensuremath{+}3.0]\\\ensuremath{+}0.08 [\ensuremath{-}0.30, \ensuremath{+}0.44]} \\
LightGBM \texttt{live\_feasible} & \shortstack{\ensuremath{+}0.4 [\ensuremath{-}2.8, \ensuremath{+}3.8]\\\ensuremath{-}0.64 [\ensuremath{-}1.48, \ensuremath{+}0.22]} & \shortstack{\ensuremath{+}0.7 [\ensuremath{-}1.8, \ensuremath{+}3.0]\\\ensuremath{-}0.23 [\ensuremath{-}0.90, \ensuremath{+}0.35]} & \shortstack{\ensuremath{+}3.2 [\ensuremath{-}0.6, \ensuremath{+}7.5]\\\ensuremath{+}0.02 [\ensuremath{-}0.66, \ensuremath{+}0.65]} & \shortstack{\ensuremath{+}0.6 [\ensuremath{-}2.4, \ensuremath{+}4.0]\\\ensuremath{-}0.12 [\ensuremath{-}0.65, \ensuremath{+}0.39]} \\
LightGBM \texttt{all\_features} & \shortstack{\ensuremath{-}1.5 [\ensuremath{-}5.4, \ensuremath{+}1.6]\\\ensuremath{-}0.24 [\ensuremath{-}0.92, \ensuremath{+}0.42]} & \shortstack{\ensuremath{-}0.1 [\ensuremath{-}2.4, \ensuremath{+}2.0]\\\ensuremath{-}0.34 [\ensuremath{-}0.76, \ensuremath{+}0.07]} & \shortstack{\ensuremath{+}0.3 [\ensuremath{-}3.1, \ensuremath{+}3.4]\\\ensuremath{+}0.10 [\ensuremath{-}0.49, \ensuremath{+}0.66]} & \shortstack{\ensuremath{+}2.3 [\ensuremath{-}0.6, \ensuremath{+}5.8]\\\ensuremath{+}0.58 [\ensuremath{-}0.14, \ensuremath{+}1.45]} \\
XGBoost \texttt{baseline} & \shortstack{\ensuremath{+}3.4 [\ensuremath{+}0.5, \ensuremath{+}6.5]\\\ensuremath{-}0.13 [\ensuremath{-}0.86, \ensuremath{+}0.50]} & \shortstack{\ensuremath{+}1.3 [\ensuremath{-}0.2, \ensuremath{+}3.1]\\\ensuremath{-}0.09 [\ensuremath{-}0.79, \ensuremath{+}0.47]} & \shortstack{\ensuremath{+}1.4 [\ensuremath{-}1.0, \ensuremath{+}3.9]\\\ensuremath{-}0.07 [\ensuremath{-}0.62, \ensuremath{+}0.48]} & \shortstack{\ensuremath{+}1.1 [\ensuremath{-}0.7, \ensuremath{+}3.1]\\\ensuremath{+}0.06 [\ensuremath{-}0.45, \ensuremath{+}0.68]} \\
XGBoost \texttt{live\_feasible} & \shortstack{\ensuremath{-}1.5 [\ensuremath{-}4.9, \ensuremath{+}1.3]\\\ensuremath{-}0.42 [\ensuremath{-}1.40, \ensuremath{+}0.50]} & \shortstack{\ensuremath{-}0.9 [\ensuremath{-}2.9, \ensuremath{+}0.9]\\\ensuremath{+}0.02 [\ensuremath{-}0.44, \ensuremath{+}0.54]} & \shortstack{\ensuremath{-}0.7 [\ensuremath{-}4.9, \ensuremath{+}3.0]\\\ensuremath{+}0.28 [\ensuremath{-}0.52, \ensuremath{+}1.17]} & \shortstack{\ensuremath{-}0.5 [\ensuremath{-}2.3, \ensuremath{+}1.3]\\\ensuremath{+}0.49 [\ensuremath{-}0.18, \ensuremath{+}1.20]} \\
XGBoost \texttt{all\_features} & \shortstack{\ensuremath{+}0.2 [\ensuremath{-}2.9, \ensuremath{+}2.8]\\\ensuremath{+}0.05 [\ensuremath{-}0.46, \ensuremath{+}0.56]} & \shortstack{\ensuremath{-}0.3 [\ensuremath{-}2.5, \ensuremath{+}1.6]\\\ensuremath{+}0.11 [\ensuremath{-}0.35, \ensuremath{+}0.60]} & \shortstack{\ensuremath{+}2.5 [\ensuremath{-}2.3, \ensuremath{+}7.8]\\\ensuremath{-}0.64 [\ensuremath{-}1.35, \ensuremath{+}0.03]} & \shortstack{\ensuremath{+}1.4 [\ensuremath{-}2.6, \ensuremath{+}5.3]\\\ensuremath{-}0.38 [\ensuremath{-}1.15, \ensuremath{+}0.28]} \\
random forest \texttt{baseline} & \shortstack{\ensuremath{-}0.2 [\ensuremath{-}2.8, \ensuremath{+}2.3]\\\ensuremath{-}0.03 [\ensuremath{-}0.51, \ensuremath{+}0.50]} & \shortstack{\ensuremath{+}0.9 [\ensuremath{-}0.8, \ensuremath{+}2.9]\\\ensuremath{-}0.07 [\ensuremath{-}0.52, \ensuremath{+}0.35]} & \shortstack{\ensuremath{+}0.8 [\ensuremath{-}1.5, \ensuremath{+}3.0]\\\ensuremath{-}0.15 [\ensuremath{-}0.63, \ensuremath{+}0.33]} & \shortstack{\ensuremath{+}1.0 [\ensuremath{-}0.7, \ensuremath{+}2.7]\\\ensuremath{+}0.11 [\ensuremath{-}0.26, \ensuremath{+}0.50]} \\
random forest \texttt{live\_feasible} & \shortstack{\ensuremath{+}0.2 [\ensuremath{-}2.0, \ensuremath{+}2.5]\\\ensuremath{-}0.25 [\ensuremath{-}0.68, \ensuremath{+}0.13]} & \shortstack{\ensuremath{-}0.4 [\ensuremath{-}2.0, \ensuremath{+}1.1]\\\ensuremath{-}0.27 [\ensuremath{-}0.63, \ensuremath{+}0.07]} & \shortstack{\ensuremath{+}2.8 [\ensuremath{-}2.4, \ensuremath{+}10.1]\\\ensuremath{-}0.57 [\ensuremath{-}1.05, \ensuremath{-}0.07]} & \shortstack{\ensuremath{-}0.7 [\ensuremath{-}2.1, \ensuremath{+}0.7]\\\ensuremath{-}0.12 [\ensuremath{-}0.63, \ensuremath{+}0.39]} \\
random forest \texttt{all\_features} & \shortstack{\ensuremath{-}1.9 [\ensuremath{-}4.5, \ensuremath{+}0.4]\\\ensuremath{+}0.33 [\ensuremath{-}0.19, \ensuremath{+}0.89]} & \shortstack{\ensuremath{-}1.1 [\ensuremath{-}2.8, \ensuremath{+}0.4]\\\ensuremath{-}0.02 [\ensuremath{-}0.55, \ensuremath{+}0.43]} & \shortstack{\ensuremath{+}1.2 [\ensuremath{-}2.3, \ensuremath{+}6.1]\\\ensuremath{+}0.30 [\ensuremath{-}0.33, \ensuremath{+}0.96]} & \shortstack{\ensuremath{-}0.4 [\ensuremath{-}1.9, \ensuremath{+}1.0]\\\ensuremath{+}0.12 [\ensuremath{-}0.20, \ensuremath{+}0.48]} \\
\bottomrule
\end{tabular}
}
\end{table}

% Source: results/linear_subsection_trees_optuna_mask/report/tuneper_oos.csv
\begin{table}[H]
\centering
\caption{How often to tune: Optuna (best of 50) with a new tuning point every 1, 5 and 25
sessions minus every 250 sessions, every path refit every session; the second column is the
QLIKE at TUNE\_PER 250. Cells as in Table~\ref{tab:app_camp_mask}.}
\label{tab:app_camp_tuneper}
\campfit{\scriptsize% AUTO-GENERATED by writeup/make_appendix_campaign_16h_tex.py -- do not edit.
% Source: results/linear_subsection_trees_optuna_mask/report/tuneper_oos.csv (comparison 'tune_per vs 250')
\begin{tabular}{lrccc}
\toprule
arm & \shortstack{QLIKE\\TUNE\_PER 250} & \shortstack{TUNE\_PER 1\\vs 250} & \shortstack{TUNE\_PER 5\\vs 250} & \shortstack{TUNE\_PER 25\\vs 250} \\
\midrule
LightGBM \texttt{baseline} & 0.1142 & \shortstack{\ensuremath{-}4.9 [\ensuremath{-}11.4, \ensuremath{+}0.4]\\\ensuremath{+}0.25 [\ensuremath{-}0.41, \ensuremath{+}0.83]} & \shortstack{\ensuremath{-}5.8 [\ensuremath{-}12.3, \ensuremath{-}0.7]\\\ensuremath{-}0.08 [\ensuremath{-}0.68, \ensuremath{+}0.49]} & \shortstack{\ensuremath{-}4.4 [\ensuremath{-}9.3, \ensuremath{-}0.3]\\\ensuremath{+}0.09 [\ensuremath{-}0.61, \ensuremath{+}0.72]} \\
LightGBM \texttt{live\_feasible} & 0.1055 & \shortstack{\ensuremath{-}2.8 [\ensuremath{-}9.2, \ensuremath{+}2.9]\\\ensuremath{-}0.55 [\ensuremath{-}1.42, \ensuremath{+}0.30]} & \shortstack{\ensuremath{-}3.4 [\ensuremath{-}9.6, \ensuremath{+}1.8]\\\ensuremath{-}0.13 [\ensuremath{-}0.66, \ensuremath{+}0.38]} & \shortstack{\ensuremath{+}1.3 [\ensuremath{-}3.2, \ensuremath{+}6.5]\\\ensuremath{-}0.10 [\ensuremath{-}0.56, \ensuremath{+}0.39]} \\
LightGBM \texttt{all\_features} & 0.1041 & \shortstack{\ensuremath{+}0.4 [\ensuremath{-}4.9, \ensuremath{+}4.7]\\\ensuremath{-}0.13 [\ensuremath{-}0.72, \ensuremath{+}0.45]} & \shortstack{\ensuremath{-}0.3 [\ensuremath{-}6.1, \ensuremath{+}4.8]\\\ensuremath{-}0.31 [\ensuremath{-}1.23, \ensuremath{+}0.62]} & \shortstack{\ensuremath{+}5.4 [\ensuremath{-}1.0, \ensuremath{+}14.5]\\\ensuremath{+}0.05 [\ensuremath{-}0.56, \ensuremath{+}0.68]} \\
XGBoost \texttt{baseline} & 0.1101 & \shortstack{\ensuremath{-}1.7 [\ensuremath{-}8.3, \ensuremath{+}4.0]\\\ensuremath{-}0.34 [\ensuremath{-}1.11, \ensuremath{+}0.33]} & \shortstack{\ensuremath{-}3.4 [\ensuremath{-}9.4, \ensuremath{+}1.3]\\\ensuremath{-}0.32 [\ensuremath{-}0.84, \ensuremath{+}0.13]} & \shortstack{\ensuremath{-}4.0 [\ensuremath{-}10.7, \ensuremath{+}1.0]\\\ensuremath{-}0.13 [\ensuremath{-}0.59, \ensuremath{+}0.31]} \\
XGBoost \texttt{live\_feasible} & 0.1024 & \shortstack{\ensuremath{-}1.6 [\ensuremath{-}7.5, \ensuremath{+}3.6]\\\ensuremath{-}0.67 [\ensuremath{-}1.62, \ensuremath{+}0.32]} & \shortstack{\ensuremath{-}1.0 [\ensuremath{-}6.8, \ensuremath{+}4.3]\\\ensuremath{-}0.62 [\ensuremath{-}1.51, \ensuremath{+}0.30]} & \shortstack{\ensuremath{+}0.1 [\ensuremath{-}5.0, \ensuremath{+}5.3]\\\ensuremath{-}0.04 [\ensuremath{-}0.55, \ensuremath{+}0.54]} \\
XGBoost \texttt{all\_features} & 0.1032 & \shortstack{\ensuremath{+}2.1 [\ensuremath{-}3.3, \ensuremath{+}6.9]\\\ensuremath{-}0.31 [\ensuremath{-}1.26, \ensuremath{+}0.59]} & \shortstack{\ensuremath{-}0.2 [\ensuremath{-}5.4, \ensuremath{+}4.5]\\\ensuremath{+}0.05 [\ensuremath{-}0.51, \ensuremath{+}0.67]} & \shortstack{\ensuremath{+}2.3 [\ensuremath{-}0.5, \ensuremath{+}5.1]\\\ensuremath{-}0.46 [\ensuremath{-}1.24, \ensuremath{+}0.27]} \\
random forest \texttt{baseline} & 0.1074 & \shortstack{\ensuremath{-}0.5 [\ensuremath{-}8.4, \ensuremath{+}5.8]\\\ensuremath{+}0.01 [\ensuremath{-}0.63, \ensuremath{+}0.74]} & \shortstack{\ensuremath{-}2.7 [\ensuremath{-}10.5, \ensuremath{+}2.7]\\\ensuremath{+}0.13 [\ensuremath{-}0.61, \ensuremath{+}0.90]} & \shortstack{\ensuremath{-}1.0 [\ensuremath{-}6.3, \ensuremath{+}3.2]\\\ensuremath{-}0.18 [\ensuremath{-}0.82, \ensuremath{+}0.41]} \\
random forest \texttt{live\_feasible} & 0.1083 & \shortstack{\ensuremath{-}7.5 [\ensuremath{-}16.7, \ensuremath{-}1.1]\\\ensuremath{-}0.06 [\ensuremath{-}0.56, \ensuremath{+}0.48]} & \shortstack{\ensuremath{-}7.1 [\ensuremath{-}16.3, \ensuremath{-}0.7]\\\ensuremath{-}0.15 [\ensuremath{-}0.57, \ensuremath{+}0.30]} & \shortstack{\ensuremath{-}4.8 [\ensuremath{-}10.5, \ensuremath{-}0.6]\\\ensuremath{-}0.62 [\ensuremath{-}1.44, \ensuremath{+}0.09]} \\
random forest \texttt{all\_features} & 0.1184 & \shortstack{\ensuremath{-}13.7 [\ensuremath{-}33.8, \ensuremath{-}1.1]\\\ensuremath{-}0.22 [\ensuremath{-}0.74, \ensuremath{+}0.27]} & \shortstack{\ensuremath{-}13.6 [\ensuremath{-}32.5, \ensuremath{-}1.5]\\\ensuremath{-}0.04 [\ensuremath{-}0.50, \ensuremath{+}0.38]} & \shortstack{\ensuremath{-}11.9 [\ensuremath{-}27.9, \ensuremath{-}1.3]\\\ensuremath{+}0.15 [\ensuremath{-}0.22, \ensuremath{+}0.57]} \\
\bottomrule
\end{tabular}
}
\end{table}

% Source: results/linear_subsection_trees_optuna_mask/report/tuneper_oos.csv
\begin{table}[H]
\centering
\caption{Optuna (best of 50, each TUNE\_PER) minus the shipped trees refit every session
(T1), the random search refit every session (RS1) and the per-bar ridge: the number of model
$\times$ input-set arms whose 95\,\% interval lies wholly below or above zero.}
\label{tab:app_camp_optuna_vs}
\campfit{\small% AUTO-GENERATED by writeup/make_appendix_campaign_16h_tex.py -- do not edit.
% Source: results/linear_subsection_trees_optuna_mask/report/tuneper_oos.csv (comparisons 'vs T1', 'vs RS1', 'vs ridge')
\begin{tabular}{llrrrrr}
\toprule
Optuna $-$ & TUNE\_PER & arms & \shortstack{QLIKE\\below 0} & \shortstack{QLIKE\\above 0} & \shortstack{Sharpe mid\\below 0} & \shortstack{Sharpe mid\\above 0} \\
\midrule
T1 & 1 & 9 & 0 & 5 & 0 & 0 \\
T1 & 5 & 9 & 0 & 5 & 0 & 0 \\
T1 & 25 & 9 & 0 & 4 & 0 & 0 \\
T1 & 250 & 9 & 0 & 5 & 0 & 0 \\
\emph{all} &  & 36 & 0 & 19 & 0 & 0 \\
\midrule
RS1 & 1 & 9 & 2 & 0 & 0 & 0 \\
RS1 & 5 & 9 & 2 & 0 & 1 & 0 \\
RS1 & 25 & 9 & 2 & 2 & 0 & 0 \\
RS1 & 250 & 9 & 2 & 1 & 0 & 0 \\
\emph{all} &  & 36 & 8 & 3 & 1 & 0 \\
\midrule
per-bar ridge & 1 & 9 & 0 & 3 & 0 & 0 \\
per-bar ridge & 5 & 9 & 0 & 3 & 1 & 0 \\
per-bar ridge & 25 & 9 & 0 & 3 & 1 & 0 \\
per-bar ridge & 250 & 9 & 0 & 5 & 0 & 0 \\
\emph{all} &  & 36 & 0 & 14 & 2 & 0 \\
\bottomrule
\end{tabular}
}
\end{table}

% Source: results/linear_subsection_trees_optuna_mask/report/picks_bounds.csv
\begin{table}[H]
\centering
\caption{Where the best-of-50 picks sit: share of picks in the outermost twentieth of each
axis at its low and high end, per input set (\texttt{baseline} / \texttt{live\_feasible} /
\texttt{all\_features}), beside the share of the random start-up draws there (the prior),
and the median pick. $^{\dagger}$A pseudo-axis: the early-stopped number of rounds at the
cap set by the learning rate.}
\label{tab:app_camp_picks}
\campfit{\scriptsize% AUTO-GENERATED by writeup/make_appendix_campaign_16h_tex.py -- do not edit.
% Source: results/linear_subsection_trees_optuna_mask/report/picks_bounds.csv
\begin{tabular}{llccccc}
\toprule
axis & search space & \shortstack{picks at low bound (\%)\\base / live / all} & \shortstack{picks at high bound (\%)\\base / live / all} & \shortstack{prior\\low (\%)} & \shortstack{prior\\high (\%)} & \shortstack{median pick\\base / live / all} \\
\midrule
\multicolumn{7}{l}{\emph{LightGBM}} \\
\texttt{num\_leaves} & [2, 128] log int & 20.6 / 7.6 / 5.7 & 2.0 / 3.3 / 3.3 & 11.1 & 4.8 & 6 / 14 / 15 \\
\texttt{min\_child\_samples} & [1, 256] log int & 19.1 / 22.7 / 14.2 & 0.7 / 0.5 / 0.8 & 18.3 & 4.2 & 5 / 5 / 12 \\
\texttt{feature\_fraction} & [0.05, 1] & 1.4 / 5.9 / 5.0 & 7.9 / 5.4 / 5.3 & 5.2 & 5.3 & 0.646 / 0.486 / 0.553 \\
\texttt{bagging\_fraction} & [0.3, 1] & 5.2 / 4.8 / 4.1 & 5.0 / 3.8 / 4.6 & 4.7 & 4.9 & 0.638 / 0.644 / 0.666 \\
\texttt{lambda\_l2} & [0.001, 10000] log & 5.4 / 5.2 / 5.6 & 0.0 / 0.0 / 0.0 & 5.5 & 5.1 & 0.633 / 0.594 / 0.637 \\
\texttt{learning\_rate} & [0.005, 0.1] log & 3.1 / 4.4 / 3.7 & 6.3 / 4.2 / 4.7 & 4.9 & 5.2 & 0.0256 / 0.02 / 0.0203 \\
\texttt{rounds\_cap}$^{\dagger}$ & early-stopped rounds $\le$ cap(lr) & --- / --- / --- & 1.1 / 0.5 / 0.6 & --- & 16.9--17.9 & 272 / 324 / 294 \\
\multicolumn{7}{l}{\emph{XGBoost}} \\
\texttt{max\_depth} & [1, 8] int & 21.4 / 8.8 / 7.4 & 6.1 / 9.5 / 7.4 & 12.2 & 12.5 & 3 / 4 / 4 \\
\texttt{min\_child\_weight} & [0.1, 256] log & 6.1 / 5.6 / 4.9 & 0.4 / 0.6 / 1.0 & 5.3 & 4.8 & 1.45 / 1.92 / 2.42 \\
\texttt{subsample} & [0.15, 1] & 6.9 / 5.4 / 3.1 & 3.9 / 4.2 / 4.8 & 5.2 & 5.3 & 0.573 / 0.598 / 0.618 \\
\texttt{colsample\_bytree} & [0.05, 1] & 0.2 / 4.5 / 3.3 & 8.3 / 4.3 / 7.0 & 4.7 & 4.9 & 0.691 / 0.538 / 0.586 \\
\texttt{reg\_lambda} & [0.0001, 10000] log & 5.5 / 6.3 / 5.9 & 0.0 / 0.0 / 0.0 & 5.5 & 5.1 & 0.101 / 0.125 / 0.144 \\
\texttt{learning\_rate} & [0.005, 0.1] log & 3.5 / 4.2 / 3.6 & 6.3 / 4.2 / 3.9 & 4.9 & 5.2 & 0.0241 / 0.0194 / 0.0213 \\
\texttt{rounds\_cap}$^{\dagger}$ & early-stopped rounds $\le$ cap(lr) & --- / --- / --- & 2.0 / 1.6 / 1.8 & --- & 18.5--20.4 & 313 / 354 / 323 \\
\multicolumn{7}{l}{\emph{random forest}} \\
\texttt{max\_features} & [0.02, 1] & 0.3 / 2.6 / 2.9 & 3.6 / 4.3 / 5.9 & 5.1 & 4.8 & 0.516 / 0.382 / 0.395 \\
\texttt{min\_samples\_leaf} & [1, 256] log int & 17.5 / 29.0 / 14.1 & 0.0 / 0.0 / 0.1 & 17.7 & 4.4 & 4 / 3 / 7 \\
\texttt{max\_depth} & \{2, 4, 8, 16, 32, None\} & 0.1 / 0.0 / 0.3 & 20.1 / 28.4 / 25.3 & 17.0 & 16.5 & 16 / 32 / 32 \\
\bottomrule
\end{tabular}
}
\end{table}

% ----------------------------------------------------------------------------
\subsubsection{The intraday-sequence LSTM}
\label{sec:app_camp_lstmi}

The per-bar LSTM reads a sequence of the last sessions' 16:00-row vectors of the per-bar
design, the HAR ladders among them. The intraday LSTM reads instead the last $N$ half-hour
bars of the near-24-hour panel ending with the 15:30 bar, one step per bar, with bar-level
inputs --- the bar's adjusted target, the adjusted exogenous values of the input set, their
availability and activity flags, the calendar columns and a half-hour clock --- so that the
network, not the ladders, does the aggregation over time. $N$ is tuned in
$\{13, 48, 96\}$ with the per-bar LSTM's grid, selection rule, window, seed average and
daily refits, and the per-window mask acts on the step columns (kept at the refits:
\texttt{baseline} \campLstmiKeptBase{} of \campLstmiStepBase{}, range
\campLstmiKeptBaseLo--\campLstmiKeptBaseHi; \texttt{live\_feasible} \campLstmiKeptLive{} of
\campLstmiStepLive{}, \campLstmiKeptLiveLo--\campLstmiKeptLiveHi; \texttt{all\_features}
\campLstmiKeptAll{} of \campLstmiStepAll{}, \campLstmiKeptAllLo--\campLstmiKeptAllHi). The
chosen $N$ moves over time: over the \campLstmiTunings{} tuning points of the three input
sets (the selection rule of record) it is 13 at \campLstmiNThirteen{}, 48 at
\campLstmiNFortyEight{} and 96 at \campLstmiNNinetySix{}.

Against the per-bar LSTM (the masked, daily-refit table), the intraday LSTM's QLIKE interval
lies below zero in \campLstmiVsLSTMQBelow{} of \campLstmiVsLSTMN{} input sets and above zero
in \campLstmiVsLSTMQAbove{}; against the per-bar ridge it lies above zero in
\campLstmiVsRidgeQAbove{} of \campLstmiVsRidgeN{}. No Sharpe interval excludes zero
(against the per-bar LSTM \campLstmiVsLSTMSBelow{} below and \campLstmiVsLSTMSAbove{}
above; against the ridge \campLstmiVsRidgeSBelow{} below and \campLstmiVsRidgeSAbove{}
above;
Table~\ref{tab:app_camp_lstmi}).

% Source: results/linear_subsection_lstm_intraday/score/lstmi_levels.csv, lstmi_pairs.csv
\begin{table}[H]
\centering
\caption{The intraday-sequence LSTM (selection rule of record) beside the per-bar LSTM and
the per-bar ridge: QLIKE and $\mathrm{sign}(s)$ Sharpe (midpoint) levels, and the paired
differences, cells as in Table~\ref{tab:app_camp_mask}.}
\label{tab:app_camp_lstmi}
\campfit{\scriptsize% AUTO-GENERATED by writeup/make_appendix_campaign_16h_tex.py -- do not edit.
% Source: results/linear_subsection_lstm_intraday/score/lstmi_levels.csv
% Source: results/linear_subsection_lstm_intraday/score/lstmi_pairs.csv
\begin{tabular}{lrrrrrrcc}
\toprule
input set & \shortstack{QLIKE\\intraday} & \shortstack{QLIKE\\per-bar LSTM} & \shortstack{QLIKE\\ridge} & \shortstack{Sharpe mid\\intraday} & \shortstack{Sharpe mid\\per-bar LSTM} & \shortstack{Sharpe mid\\ridge} & \shortstack{intraday $-$\\per-bar LSTM} & \shortstack{intraday $-$\\ridge} \\
\midrule
\texttt{baseline} & 0.1073 & 0.1049 & 0.0982 & 0.81 & 1.26 & 1.22 & \shortstack{\ensuremath{+}2.2 [\ensuremath{-}3.4, \ensuremath{+}9.0]\\\ensuremath{-}0.45 [\ensuremath{-}1.07, \ensuremath{+}0.20]} & \shortstack{\ensuremath{+}9.2 [\ensuremath{+}2.5, \ensuremath{+}17.2]\\\ensuremath{-}0.41 [\ensuremath{-}1.18, \ensuremath{+}0.41]} \\
\texttt{live\_feasible} & 0.1103 & 0.1193 & 0.1005 & 1.53 & 0.91 & 1.90 & \shortstack{\ensuremath{-}7.5 [\ensuremath{-}20.9, \ensuremath{+}4.7]\\\ensuremath{+}0.62 [\ensuremath{-}0.26, \ensuremath{+}1.54]} & \shortstack{\ensuremath{+}9.7 [\ensuremath{+}0.1, \ensuremath{+}20.0]\\\ensuremath{-}0.37 [\ensuremath{-}1.33, \ensuremath{+}0.60]} \\
\texttt{all\_features} & 0.1119 & 0.1305 & 0.1004 & 0.97 & 1.33 & 1.66 & \shortstack{\ensuremath{-}14.3 [\ensuremath{-}26.9, \ensuremath{-}2.0]\\\ensuremath{-}0.37 [\ensuremath{-}1.42, \ensuremath{+}0.76]} & \shortstack{\ensuremath{+}11.5 [\ensuremath{+}2.2, \ensuremath{+}22.5]\\\ensuremath{-}0.70 [\ensuremath{-}1.74, \ensuremath{+}0.35]} \\
\bottomrule
\end{tabular}
}
\end{table}

% ----------------------------------------------------------------------------
\subsubsection{Reproducibility across CPUs}
\label{sec:app_camp_repro}

The tree and LSTM numbers depend on the CPU's vector class. The executor builds the design
matrix $X$ with numpy's runtime-dispatched loops, and $X$ comes out as one of two byte
strings: one on every machine with AVX-512 and another on every machine without it
(\campXHosts{} machines probed on two clusters). The two differ in \campXColsLive{} of
\campPLive{} \texttt{live\_feasible} columns and \campXColsAll{} of \campPAll{}
\texttt{all\_features} columns, by at most \campXMaxAbs{} in absolute value; the target and
the kept-column sets are the same on every machine. Within a class the Optuna trials and
forecasts reproduce bit for bit across clusters (\campSameClassTrials{} trials,
\campSameClassFcst{} forecasts). Across classes the fits part: from the same stage-1 records
the random forest's forecasts are identical in \campXFcstRf{}, XGBoost's in \campXFcstXgb{}
(largest relative difference \campXFcstMaxXgb) and LightGBM's in \campXFcstLgbm{}
(\campXFcstMaxLgbm); re-running \campXcPoints{} tuning points on the other class changes
LightGBM's best-of-50 pick at \campXcPickLgbm{} of them (XGBoost \campXcPickXgb{}, random
forest \campXcPickRf{}), while a re-run on the same class is bit-identical at all
\campXcSame{}. A whole masked LightGBM \texttt{live\_feasible} run on the other class moves
the forecast by up to \campClassRelTen\,\% (T10) and \campClassRelOne\,\% (T1), QLIKE by
\campClassQTen\,\% \campClassQTenCI{} and \campClassQOne\,\% \campClassQOneCI{}, and the
Sharpe ratio by \campClassSTen{} \campClassSTenCI{} and \campClassSOne{} \campClassSOneCI{}
(Table~\ref{tab:app_camp_repro}).

So every canonical chunk ran on one CPU class: the census finds \campHChunkClasses{} class
among the masked ladder's \campHChunks{} chunks and \campCChunkClasses{} among Optuna's
\campCChunks{} chunk-arms (the masked ladder's \campHGateChunks{} class-gate chunks ran on
the other class on purpose), and the intraday LSTM, whose first
attempt on mixed nodes was stopped by its tuning cache's fingerprint, was re-run on one class
(the mixed attempt agrees on \campLstmiXTune{} tuning points and \campLstmiXChunk{} refit
chunks). Optuna's first pass had landed on several classes; it is kept as the class-mixed
run beside the one-class run: \campSMDiffer{} of its \campSMN{} forecast tables differ, with
\campSMQBelow{} QLIKE intervals below zero and \campSMQAbove{} above (one class minus mixed)
and \campSMSharpeExcl{} Sharpe intervals excluding zero.

% Sources: see the generated file (results/linear_subsection_trees_optuna/crosscluster/,
% results/linear_subsection_trees_optuna_mask/xclass*/, results/linear_subsection_trees_mask/class/,
% results/trees_mask_1600/mask_pairs.csv, results/linear_subsection_lstm_intraday/gates/,
% results/linear_subsection_trees_optuna_mask/report/single_vs_mixed.csv)
\begin{table}[H]
\centering
\caption{CPU-class checks: what was compared and what reproduced. ``Bit-identical $a/b$'':
$a$ of $b$ items equal bit for bit; relative differences are of the forecast; QLIKE and
Sharpe differences (other class minus the fleet's class, in \% of QLIKE and Sharpe units)
with 95\,\% paired intervals.}
\label{tab:app_camp_repro}
\campfit{\small% AUTO-GENERATED by writeup/make_appendix_campaign_16h_tex.py -- do not edit.
% Source: results/linear_subsection_trees_optuna/crosscluster/arch/design_diffs.csv
% Source: results/linear_subsection_trees_optuna/crosscluster/arch/probe_hosts.csv
% Source: results/linear_subsection_trees_optuna/crosscluster/mask/qualify_h2x512_vs_carc_xeon.csv (+ _h2x512r_, s2class/)
% Source: results/linear_subsection_trees_optuna/crosscluster/mask/qualify_carc_xeon_vs_carc_epyc.csv (+ s2class/)
% Source: results/linear_subsection_trees_optuna_mask/xclass/xclass_gate.csv (+ xclass_epyc/)
% Source: results/linear_subsection_trees_mask/class/class_arms.csv, results/trees_mask_1600/mask_pairs.csv (family 'CPU class')
% Source: results/linear_subsection_lstm_intraday/gates/cross_class_mixed_vs_pinned.csv
% Source: results/linear_subsection_trees_optuna_mask/report/single_vs_mixed.csv
\begin{tabular}{p{0.30\textwidth}p{0.64\textwidth}}
\toprule
check & result \\
\midrule
design $X$, machines of the two vector classes & \texttt{live\_feasible}: \campXColsLive{} of \campPLive{} columns differ; \texttt{all\_features}: \campXColsAll{} of \campPAll{}; largest absolute difference \campXMaxAbs{}; target and kept-column sets identical \\
Optuna, same class on two clusters & stage-1 trials bit-identical \campSameClassTrials{} (replicate \campSameClassTrialsRep{}); stage-2 forecasts \campSameClassFcst{} \\
Optuna, across classes on one cluster (same inputs) & stage-1 trials bit-identical: LightGBM \campXTrialsLgbm{}, XGBoost \campXTrialsXgb{}, random forest \campXTrialsRf{}; stage-2 forecasts from the same records: \campXFcstLgbm{} (max rel. \campXFcstMaxLgbm{}), \campXFcstXgb{} (\campXFcstMaxXgb{}), \campXFcstRf{} \\
Optuna, first \campXcPoints{} tuning points re-run on the other class & best-of-50 pick changes: LightGBM \campXcPickLgbm{}, XGBoost \campXcPickXgb{}, random forest \campXcPickRf{}; XGBoost forecast max rel. \campXcFcstXgb{}; same-class re-run bit-identical at \campXcSame{} of \campXcPoints{} points for all three \\
masked LightGBM \texttt{live\_feasible}, whole run on the other class & T10 / T1: max rel. forecast difference \campClassRelTen{} / \campClassRelOne{}\,\%; QLIKE \campClassQTen{} \campClassQTenCI{} / \campClassQOne{} \campClassQOneCI{}\,\%; $\Delta$Sharpe mid \campClassSTen{} \campClassSTenCI{} / \campClassSOne{} \campClassSOneCI{} \\
intraday LSTM, mixed-class attempt vs one class & tuning points bit-identical \campLstmiXTune{}; refit chunks \campLstmiXChunk{} (the others ran on the other class) \\
Optuna, one class vs class-mixed first pass & \campSMDiffer{} of \campSMN{} tables differ (max rel. LightGBM \campSMMaxLgbm{}, XGBoost \campSMMaxXgb{}, random forest \campSMMaxRf{}); QLIKE intervals below / above 0: \campSMQBelow{} / \campSMQAbove{}; Sharpe mid intervals excluding 0: \campSMSharpeExcl{} \\
\bottomrule
\end{tabular}
}
\end{table}

% ----------------------------------------------------------------------------
\subsubsection{The master table before and after}
\label{sec:app_camp_master}

The master table of the closing strategy was rebuilt from the new tables, on the same
\campMTHeadDays{} days and with the same resampled days for every interval. The headline
forecast, the per-bar ridge on \texttt{live\_feasible}, is unchanged: QLIKE
\campMTHeadQBefore{} before and \campMTHeadQ{} after, $\mathrm{sign}(s)$ Sharpe
\campMTHeadS{} at the midpoint and \campMTHeadSX{} crossed, positions changed on
\campMTHeadPos{} of \campMTHeadDays{} days (largest relative change of its recalibrated
forecast \campMTHeadMaxRel). Table~\ref{tab:app_camp_master} summarises the
\campMTKeysBoth{} forecasts present before and after, by family; \campMTKeysNew{} forecasts
are new (the daily-refit and random-search-daily trees, the Optuna rungs and the intraday
LSTM). The own Sharpe ratio of \campMTSigN{} forecasts moved with the interval excluding
zero: \campMTSigWho.

Against the headline, of the \campMTNOther{} other forecasts \campMTBeatMid{} have a
Sharpe-difference interval wholly above zero at the midpoint and \campMTBeatCrossed{}
crossed, and \campMTWorseMid{} have one wholly below zero at the midpoint. On forecast
accuracy \campMTQBetter{} forecasts have a QLIKE interval against it wholly below zero, from
the family \campMTQBetterFamily{}, none of them with a Sharpe interval above zero:
\campMTQBetterWho{}. \campMTQWorse{} have a QLIKE interval wholly above zero.

% Source: results/close_master_table/before_after.csv
\begin{table}[H]
\centering
\caption{The master table before and after the campaign, per family of the forecasts present
in both: what changed, the QLIKE change in \% and the $\mathrm{sign}(s)$ Sharpe change
(midpoint), median and range over the family, the number of paired Sharpe-change intervals
wholly above / below zero, and the trade positions changed (of \campMTHeadDays{} days).}
\label{tab:app_camp_master}
\campfit{\scriptsize% AUTO-GENERATED by writeup/make_appendix_campaign_16h_tex.py -- do not edit.
% Source: results/close_master_table/before_after.csv (status 'both')
\begin{tabular}{lrp{0.22\textwidth}cccc}
\toprule
family & keys & what changed & \shortstack{QLIKE change (\%)\\median [min, max]} & \shortstack{$\Delta$Sharpe mid\\median [min, max]} & \shortstack{$\Delta$Sharpe interval\\above / below 0} & \shortstack{positions changed\\median [max]} \\
\midrule
paper & 9 & nothing & 0.00 [0.00, 0.00] & 0.00 [0.00, 0.00] & 0 / 0 & 0 [0] \\
per-bar linear & 10 & design de-dup & 0.00 [0.00, \ensuremath{+}0.05] & 0.00 [0.00, 0.00] & 0 / 0 & 0 [0] \\
pooled twin & 9 & nothing & 0.00 [0.00, 0.00] & 0.00 [0.00, 0.00] & 0 / 0 & 0 [0] \\
VIX-only family & 21 & design de-dup & 0.00 [\ensuremath{-}0.04, \ensuremath{+}0.09] & 0.00 [\ensuremath{-}0.05, \ensuremath{+}0.06] & 0 / 0 & 0 [2] \\
per-bar tree (untuned) & 9 & design de-dup + per-window mask + CPU-class pinning & \ensuremath{+}0.63 [\ensuremath{+}0.26, \ensuremath{+}1.64] & \ensuremath{-}0.10 [\ensuremath{-}0.43, \ensuremath{+}0.49] & 0 / 2 & 26 [49] \\
per-bar tree (tuned) & 18 & design de-dup + per-window mask + CPU-class pinning & \ensuremath{+}0.94 [\ensuremath{-}5.79, \ensuremath{+}16.14] & \ensuremath{+}0.07 [\ensuremath{-}0.27, \ensuremath{+}0.48] & 0 / 0 & 48 [80] \\
LSTM & 6 & design de-dup + per-window mask + refit cadence 10 $\to$ 1 + CPU-class pinning & \ensuremath{-}4.40 [\ensuremath{-}17.46, \ensuremath{-}1.68] & \ensuremath{+}0.40 [\ensuremath{-}0.16, \ensuremath{+}0.80] & 1 / 0 & 128 [172] \\
direct rest-of-day at 15:30 (check) & 9 & design de-dup + CPU-architecture pinning to the per-bar twin & 0.00 [0.00, 0.00] & 0.00 [\ensuremath{-}0.01, 0.00] & 0 / 0 & 0 [1] \\
implied-vol representations & 15 & design de-dup & 0.00 [0.00, \ensuremath{+}0.06] & 0.00 [0.00, 0.00] & 0 / 0 & 0 [0] \\
HAR-ladder variants & 9 & design de-dup & 0.00 [\ensuremath{-}0.01, \ensuremath{+}0.01] & 0.00 [0.00, 0.00] & 0 / 0 & 0 [1] \\
chain-period buckets & 21 & design de-dup & 0.00 [0.00, 0.00] & 0.00 [0.00, 0.00] & 0 / 0 & 0 [0] \\
\bottomrule
\end{tabular}
}
\end{table}

% ----------------------------------------------------------------------------
\subsubsection{Feature importance and density on the new design}
\label{sec:app_camp_featimp}

The importance study of the 16:00 forecast (MDI, split count, permutation, TreeSHAP) and the
dense-versus-sparse study were re-run on the new design with every tree fit masked, with the
first pass's refits, held-out tails, permutation scheme and seeds. In the trees
(Table~\ref{tab:app_camp_featimp}) the \texttt{har\_ma\_*} ladder still leads
\campFILead{} of \campFILeadN{} input-set $\times$ model $\times$ measure rows at the series
level. The table gives its MDI, TreeSHAP and permutation shares before and after, beside
the part of the old ladder's MDI that sat on the \texttt{har\_ma\_k\_x\_close} copies
(column five).
The rank correlation of the series rankings, old against new, has median
\campFISpearmanTrees{} over the tree rows (lowest \campFISpearmanTreesMin{}). No interval is
computed for these share changes; the yardstick is the ridge, which solves the same problem
in both runs, and whose permutation importance of the ladder moves by at most
\campFIRidgeNoise{} points of the tail loss between the two sets of permutation draws.

In the density table (Table~\ref{tab:app_camp_dvs}) the linear forecasts are the first
pass's to \campDvsLinMaxRel{}, so their effective number of inputs $N_{\mathrm{eff}}$ does
not move; the trees' $N_{\mathrm{eff}}$ is given for the first pass, the de-duplicated design
without the mask and with it, beside the share of the first pass's TreeSHAP that sat on the
session-edge columns (last column). \campDvsClaimsSame{} of the
first pass's \campDvsClaims{} statements read the same on the new runs. LightGBM on all
columns, new minus first pass: QLIKE \campDvsLgbmQLive{} \campDvsLgbmQLiveCI{} and Sharpe
\campDvsLgbmSLive{} \campDvsLgbmSLiveCI{} (\texttt{live\_feasible}); QLIKE
\campDvsLgbmQAll{} \campDvsLgbmQAllCI{} and Sharpe \campDvsLgbmSAll{} \campDvsLgbmSAllCI{}
(\texttt{all\_features}).

% Source: results/feature_importance_1530_dedup/before_after_har_ma.csv, before_after_har_ma_columns.csv
\begin{table}[H]
\centering
\caption{The \texttt{har\_ma\_*} ladder's importance in the trees, first pass (old design,
unmasked) to the new design with the mask: MDI and TreeSHAP shares, permutation (P2,
QLIKE) in \% of the model's tail loss, the part of the old ladder's MDI that sat on the
\texttt{\_x\_close} copies, and the ladder's MDI rank among the series.}
\label{tab:app_camp_featimp}
\campfit{\scriptsize% AUTO-GENERATED by writeup/make_appendix_campaign_16h_tex.py -- do not edit.
% Source: results/feature_importance_1530_dedup/before_after_har_ma.csv
% Source: results/feature_importance_1530_dedup/before_after_har_ma_columns.csv (measure 'mdi')
\begin{tabular}{llccccc}
\toprule
input set & model & \shortstack{MDI share (\%)\\old $\to$ new} & \shortstack{TreeSHAP share (\%)\\old $\to$ new} & \shortstack{perm.~P2 QLIKE\\(\% of tail loss)} & \shortstack{old \texttt{\_x\_close} share\\of the ladder's MDI (\%)} & \shortstack{MDI rank\\old $\to$ new} \\
\midrule
\texttt{live\_feasible} & LightGBM & 62 $\to$ 57 & 56 $\to$ 52 & \ensuremath{+}268 $\to$ \ensuremath{+}219 & 29 & 1 $\to$ 1 \\
\texttt{live\_feasible} & XGBoost & 73 $\to$ 69 & 63 $\to$ 59 & \ensuremath{+}252 $\to$ \ensuremath{+}210 & 22 & 1 $\to$ 1 \\
\texttt{live\_feasible} & random forest & 67 $\to$ 67 & 64 $\to$ 64 & \ensuremath{+}331 $\to$ \ensuremath{+}320 & 49 & 1 $\to$ 1 \\
\texttt{all\_features} & LightGBM & 57 $\to$ 52 & 47 $\to$ 44 & \ensuremath{+}237 $\to$ \ensuremath{+}200 & 22 & 1 $\to$ 1 \\
\texttt{all\_features} & XGBoost & 69 $\to$ 64 & 55 $\to$ 51 & \ensuremath{+}224 $\to$ \ensuremath{+}187 & 22 & 1 $\to$ 1 \\
\texttt{all\_features} & random forest & 62 $\to$ 62 & 58 $\to$ 58 & \ensuremath{+}290 $\to$ \ensuremath{+}283 & 54 & 1 $\to$ 1 \\
\bottomrule
\end{tabular}
}
\end{table}

% Source: results/dense_vs_sparse_dedup/before_after_density.csv, before_after_session_edge.csv
\begin{table}[H]
\centering
\caption{Signal density before and after: effective number of inputs
$N_{\mathrm{eff}}=1/\sum_j s_j^2$ of each forecast's mean absolute contributions (first pass;
new design without the mask, trees only; new design with the mask), the fewest columns
reproducing 95\,\% of the input-driven forecast's variance ($k_{95}$), the columns with any
weight, and the first pass's share of the trees' mean $|$TreeSHAP$|$ on the session-edge
columns.}
\label{tab:app_camp_dvs}
\campfit{\scriptsize% AUTO-GENERATED by writeup/make_appendix_campaign_16h_tex.py -- do not edit.
% Source: results/dense_vs_sparse_dedup/before_after_density.csv (level 'column', sample 'all forecasts')
% Source: results/dense_vs_sparse_dedup/before_after_session_edge.csv
\begin{tabular}{llrrrccc}
\toprule
input set & model & \shortstack{$N_{\mathrm{eff}}$\\first pass} & \shortstack{$N_{\mathrm{eff}}$\\de-dup, no mask} & \shortstack{$N_{\mathrm{eff}}$\\de-dup + mask} & \shortstack{$k_{95}$\\old $\to$ new} & \shortstack{columns with\\weight} & \shortstack{first pass: TreeSHAP\\on session-edge (\%)} \\
\midrule
\texttt{live\_feasible} & ridge & 40.7 &  & 40.7 & 22 $\to$ 22 & 138 $\to$ 138 &  \\
\texttt{live\_feasible} & lasso & 10.8 &  & 10.8 & 7 $\to$ 7 & 86 $\to$ 86 &  \\
\texttt{live\_feasible} & elastic net & 7.7 &  & 7.7 & 5 $\to$ 5 & 52 $\to$ 52 &  \\
\texttt{live\_feasible} & LightGBM & 19.1 & 13.9 & 14.1 & 12 $\to$ 10 & 150 $\to$ 139 & 13.0 \\
\texttt{live\_feasible} & XGBoost & 12.7 & 9.2 & 9.7 & 10 $\to$ 8 & 146 $\to$ 138 & 12.8 \\
\texttt{live\_feasible} & random forest & 13.2 & 6.8 & 6.8 & 6 $\to$ 4 & 174 $\to$ 147 & 29.1 \\
\texttt{all\_features} & ridge & 124.2 &  & 124.2 & 64 $\to$ 64 & 379 $\to$ 379 &  \\
\texttt{all\_features} & lasso & 21.3 &  & 21.3 & 6 $\to$ 6 & 204 $\to$ 204 &  \\
\texttt{all\_features} & elastic net & 39.4 &  & 39.4 & 20 $\to$ 20 & 238 $\to$ 238 &  \\
\texttt{all\_features} & LightGBM & 25.8 & 20.4 & 20.4 & 16 $\to$ 12 & 411 $\to$ 367 & 8.6 \\
\texttt{all\_features} & XGBoost & 16.4 & 11.9 & 13.0 & 15 $\to$ 11 & 392 $\to$ 358 & 11.9 \\
\texttt{all\_features} & random forest & 16.4 & 8.5 & 8.4 & 9 $\to$ 6 & 566 $\to$ 388 & 27.6 \\
\bottomrule
\end{tabular}
}
\end{table}

% ----------------------------------------------------------------------------
\subsubsection{Compute used}
\label{sec:app_camp_compute}

Table~\ref{tab:app_camp_compute} lists the compute each stage reported. Cluster A is the
batch cluster that ran the tree and LSTM fleets (Slurm); cluster B the one that ran the
linear and rest-of-day arms (SGE). Allocated CPU-hours count every core a job held for its
whole run; used CPU-hours the busy share. The largest allocation is the masked Optuna run's
(\campComputeOptunaAlloc{} allocated CPU-hours, at up to \campComputeOptunaPeak{} CPUs at
once).
% 2026-09-30 L4: left out of the table, not sourced -- the de-duplicated linear arms
% (cluster B, fleet 14970159) and the LSTM on the de-duplicated design (cluster A,
% 12479507-09) of agent A reported no CPU-hours; the master table (J) and the
% trading-test / scorer fill-ins (K) ran locally and reported none.

% Source: see the generated file (cluster_usage*.csv of each stage, E's SUMMARY.md accounting,
% F's commit message 4fa0697, G's numbers.json, D's line in writeup/PROGRESS_2026-09-29.md).
% Cluster A = CARC (usc-discovery), cluster B = Hoffman2.
\begin{table}[H]
\centering
\caption{Compute used by the campaign's stages, as each stage recorded it (--- = not
reported). For cluster B, slot-hours are wall-clock hours of the job slots.}
\label{tab:app_camp_compute}
\campfit{\small% AUTO-GENERATED by writeup/make_appendix_campaign_16h_tex.py -- do not edit.
% Source: results/trees_cadence_1600/cluster_usage.csv (checked against its SUMMARY.md)
% Source: results/linear_subsection_trees_optuna_mask/cluster_usage_unmasked_run.csv
% Source: results/trees_mask_1600/cluster_usage.csv (checked against its SUMMARY.md)
% Source: results/linear_subsection_trees_optuna_mask/cluster_usage.csv, results/linear_subsection_trees_optuna_mask/cluster_usage_single_class_rerun.csv
% Source: results/linear_subsection_lstm_intraday/score/cluster_usage.csv (checked against its SUMMARY.md)
% Source: writeup/PROGRESS_2026-09-29.md (agent log, agent D's line)
% Source: results/linear_subsection_restofday_dedup/SUMMARY.md (its accounting lines, from qacct)
% Source: git log -1 4fa0697 (agent F's commit message)
% Source: results/dense_vs_sparse_dedup/numbers.json (timing)
\begin{tabular}{llrrrl}
\toprule
stage & cluster & \shortstack{allocated\\CPU-h} & \shortstack{used\\CPU-h} & \shortstack{peak\\CPUs} & wall-clock (first start--last end) \\
\midrule
trees T10 / T1 / RS10 / RS1, no mask & A & 107.3 & 54.0 & 360 & 09-29 19:39--09-29 21:02 (1 h 23 min) \\
Optuna, no mask (stopped; not used) & A & 928.3 & --- & 1260 & 09-29 20:15--09-29 22:28 (2 h 13 min) \\
trees and LSTM with the mask, one CPU class & A & 186.4 & 102.5 & 408 & 09-29 22:09--09-30 00:50 (2 h 41 min) \\
Optuna with the mask, all jobs & A & 1422.6 & --- & 900 & 09-29 22:07--09-30 04:05 (5 h 58 min) \\
\quad of which the one-class re-run & A & 274.3 & --- & 860 & 09-30 02:44--09-30 04:05 (1 h 21 min) \\
intraday-sequence LSTM & A & 197.3 & 159.2 & 340 & 09-29 23:31--09-30 03:46 (4 h 15 min) \\
cross-cluster check (Optuna) & A & 18.2 & --- & 73 & not reported \\
cross-cluster check (Optuna) & B & 29.1 slot-h & 6.7 & 58 slots & not reported \\
rest-of-day arms, all entry clocks & B & 2.06 slot-h & 1.58 & not reported & 21:47--21:59 (11.2 min) \\
\quad attribution control & B & 0.21 slot-h & 0.18 & not reported & 22:35--22:50 (14.7 min) \\
feature importance re-run & A & 10.8 & 7.1 & 45 & 21:50--22:11 \\
dense-vs-sparse re-run & A & 4.1 & --- & 20 & 18 min of job time \\
\quad its local part & local & --- & 4.6 & 4 processes & 69 min \\
\bottomrule
\end{tabular}
}
\end{table}
